\documentclass[A4paper,12pt,twoside]{article}
\usepackage{geometry}
\geometry{hmargin=0.8in,vmargin=0.9in}
%\addtolength{\hoffset}{1.cm}
%\addtolength{\textwidth}{3.cm}
%\addtolength{\voffset}{-1.5cm}
%\addtolength{\textheight}{3.0cm}
%\geometry{letterpaper}
\usepackage{fancyhdr}
\usepackage{graphicx}
%\usepackage[colorinlistoftodos]{todonotes}
\usepackage{amssymb}
\usepackage{amsmath}
\usepackage{amsfonts}
%\usepackage{amsthm}
\usepackage{mathrsfs}
\usepackage{mathtools}
\usepackage{bm}
\usepackage{color}
\usepackage{setspace}
\usepackage{exscale}
\usepackage{relsize}
\usepackage{epstopdf}
\usepackage{float}
\usepackage{cite}
\usepackage{nicefrac}
\usepackage{setspace}
\usepackage{algorithm,algorithmic}
\usepackage{tikz}
\usetikzlibrary{positioning}
\usepackage{xcolor}
\usepackage{rotating}
\usepackage{enumerate}
\usepackage{hyperref}
\usepackage{authblk}
\usepackage{soul}
%\usepackage{indentfirst}
%\usepackage{ulem}
%\usepackage{blindtext}
\usepackage{textcomp}
\usepackage[title]{appendix}

\usepackage{cleveref}
\usepackage{refcheck}

\usepackage{theorem}
\usepackage{empheq}
\DeclareGraphicsRule{.pstex}{eps}{*}{}

%%% PACKAGES
\usepackage{booktabs,multirow} % for much better looking tables
\usepackage{array} % for better arrays (eg matrices) in maths
\usepackage{paralist} % very flexible & customisable lists (eg. enumerate/itemize, etc.)
\usepackage{verbatim} % adds environment for commenting out blocks of text & for better verbatim
\usepackage{subfigure} % make it possible to include more than one captioned figure/table in a single float
% These packages are all incorporated in the memoir class to one degree or another...

\usepackage[normalem]{ulem}

%%% HEADERS & FOOTERS
\pagestyle{fancy}
\fancyhead{} % clear all header fields
\fancyhead[LE]{\thepage}
\fancyhead[RO]{\thepage}
\fancyhead[CE]{\sc A. Chertock, S. Karni, A. Kurganov \& L. Micalizzi}
\fancyhead[CO]{\sc A New AP Method}
\fancyfoot{} % clear all footer fields
%\fancyfoot[C]{\thepage}
\renewcommand{\headrulewidth}{0.0pt}
\renewcommand{\footrulewidth}{0.0pt}
\addtolength{\headheight}{3pt}
%% END Article customise


\newtheorem{lemma}{Lemma}[section]
\newtheorem{proposition}{Proposition}[section]
\newtheorem{theorem}{Theorem}[section]
\newtheorem{remark}{Remark}[section]
\newtheorem{corollary}{Corollary}[section]
\newtheorem{definition}{Definition}[section]
%\newtheorem{algorithm}{Algorithm}[section]

%{\theorembodyfont{\rmfamily}\newtheorem{remark}{Remark}[section]}
%{\theorembodyfont{\rmfamily}\newtheorem{algorithm}{Algorithm}[section]}

\newenvironment{proof}{{\flushleft \bf Proof:}}{}
\newenvironment{proof1}{{\flushleft \bf Proof of (ii):}}{}
\newenvironment{acknowledgment}{{\flushleft \bf Acknowledgment:}}{}

\renewcommand{\theequation}{\arabic{section}.\arabic{equation}}
\renewcommand{\thefigure}{\arabic{section}.\arabic{figure}}
\renewcommand{\thetable}{\arabic{section}.\arabic{table}}

\newenvironment{note}{{\flushleft \bf Note:}}{}
\renewcommand{\arraystretch}{1.28}
\setlength{\tabcolsep}{1.25mm}
%
%\allowdisplaybreaks
\allowdisplaybreaks[1]
%
\numberwithin{equation}{section}
\numberwithin{figure}{section}
\numberwithin{table}{section}
%

\usepackage{ae,aecompl}
\def\hf {\frac{1}{2}}
\def\thf {\frac{3}{2}}
\newcommand{\ph}{+\frac{1}{2}}
\newcommand{\mh}{-\frac{1}{2}}
\newcommand{\pth}{+\frac{3}{2}}
\newcommand{\mth}{-\frac{3}{2}}
\def\mom {\mbox{\boldmath$\omega$}}
\def\mo {\bm{0}}
\def\muu {\bm{u}}
\def\mw {\bm{w}}
\def\mz {\bm{z}}
\def\Chi {{\Large{\mbox{$\chi$}}}}
\newcommand{\hb}{\highlightboxed}
\def\bnabla{\bm{\nabla}}
\newcommand{\ds}{\displaystyle}
\newcommand{\cp}{\concept}
\newcommand{\dtau}{\Delta \tau}
\newcommand{\dxi}{\Delta \xi}
\newcommand{\deta}{\Delta \eta}
\newcommand{\dt}{\Delta t}
\newcommand{\dx}{\Delta x}
\newcommand{\dy}{\Delta y}
\newcommand{\kph}{{k+\frac{1}{2}}}
\newcommand{\kmh}{{k-\frac{1}{2}}}
\newcommand{\jph}{{j+\frac{1}{2}}}
\newcommand{\jmh}{{j-\frac{1}{2}}}
\newcommand{\nmh}{{n-\frac{1}{2}}}
\newcommand{\nph}{{n+\frac{1}{2}}}
\newcommand{\jpmh}{j\pm\frac{1}{2}}
\newcommand{\dstl}{\displaystyle}
\newcommand{\bu}{\bm{w}}
\newcommand{\bH}{\bm{H}}
\newcommand{\Jph}{{J+\frac{1}{2}}}
\newcommand{\jphh}{{j+\frac{3}{2}}}
\newcommand{\jmhh}{{j-\frac{3}{2}}}
\newcommand{\Jphh}{{J+\frac{3}{2}}}
\newcommand{\Jphhh}{{J+\frac{5}{2}}}
\newcommand{\Jmh}{{J-\frac{1}{2}}}
\newcommand{\mq}{\bm{u}}
\newcommand{\mqq}{\bm{q}}
\newcommand{\mf}{\bm{f}}
\newcommand{\mF}{\bm{F}}
\newcommand{\mS}{\bm{S}}
\newcommand{\mg}{\bm{g}}
\newcommand{\mG}{\bm{G}}
\newcommand{\ms}{\bm{s}}
\newcommand{\mH}{\bm{H}}
\newcommand{\ml}{\bm{{\cal P}}}
\newcommand{\mU}{\bm{U}}
\newcommand{\mR}{\bm{R}}
\newcommand{\mL}{\bm{L}}
\newcommand{\mc}{\bm{c}}
\newcommand{\bV}{\bm{V}}
\newcommand{\Ups}{\bm{\Upsilon}}
\newcommand{\cS}{\mathcal{S}}
\newcommand{\cH}{\mathcal{H}}
\newcommand{\cG}{\mathcal{G}}
\newcommand{\cP}{\mathcal{P}}
\newcommand{\ma}{\mbox{\boldmath$\alpha$}}
\newcommand{\mph}{\mbox{\boldmath$\varphi$}}
\newcommand{\ve}{\varepsilon}

\newcommand{\xjt}{{\tilde{x}_j}}
\newcommand{\xjpt}{{\tilde{x}_{j+1}}}
\newcommand{\xjpht}{{\tilde{x}_{j+\frac{1}{2}}}}
%subindex
\newcommand{\jke}{{_{j,k}^E}}

%\xi
\newcommand{\xijph}{{\xi_{j+\frac{1}{2}}}}
\newcommand{\etakph}{{\eta_{k+\frac{1}{2}}}}

%x
\newcommand{\xjphkph}{{x_{j+\frac{1}{2},k+\frac{1}{2}}}}
\newcommand{\xjphkmh}{{x_{j+\frac{1}{2},k-\frac{1}{2}}}}
\newcommand{\xjmhkph}{{x_{j-\frac{1}{2},k+\frac{1}{2}}}}
\newcommand{\xjmhkmh}{{x_{j-\frac{1}{2},k-\frac{1}{2}}}}
\newcommand{\xjk}{{x_{j,k}}}
\newcommand{\xjkp}{{x_{j,k}^p}}
%--
\newcommand{\xj}{{x_j}}
\newcommand{\xjp}{{x_{j+1}}}
\newcommand{\xjm}{{x_{j-1}}}
\newcommand{\xjph}{{x_{j+\frac{1}{2}}}}
\newcommand{\xjmh}{{x_{j-\frac{1}{2}}}}

%y
\newcommand{\yjphkph}{{y_{j+\frac{1}{2},k+\frac{1}{2}}}}
\newcommand{\yjphkmh}{{y_{j+\frac{1}{2},k-\frac{1}{2}}}}
\newcommand{\yjmhkph}{{y_{j-\frac{1}{2},k+\frac{1}{2}}}}
\newcommand{\yjmhkmh}{{y_{j-\frac{1}{2},k-\frac{1}{2}}}}
\newcommand{\yjk}{{y_{j,k}}}
\newcommand{\yjkp}{{y_{j,k}^p}}
%--
\newcommand{\yk}{{y_k}}
\newcommand{\ykp}{{y_{k+1}}}
\newcommand{\ykm}{{y_{k-1}}}
\newcommand{\ykph}{{y_{k+\frac{1}{2}}}}
\newcommand{\ykmh}{{y_{k-\frac{1}{2}}}}

%cells
\newcommand{\Tjk}{{T_{j,k}}}
\newcommand{\Tjkp}{{T_{j,k}^p}}
\newcommand{\Tjke}{{T_{j,k}^E}}
\newcommand{\Tjkw}{{T_{j,k}^W}}
\newcommand{\Tjks}{{T_{j,k}^S}}
\newcommand{\Tjkn}{{T_{j,k}^N}}
%area
\newcommand{\Ajk}{{A_{j,k}}}
\newcommand{\Ajkp}{{A_{j,k}^p}}
\newcommand{\Ajke}{{A_{j,k}^E}}
\newcommand{\Ajkw}{{A_{j,k}^W}}
\newcommand{\Ajks}{{A_{j,k}^S}}
\newcommand{\Ajkn}{{A_{j,k}^N}}
%center of mass
\newcommand{\Cjk}{{C_{j,k}}}
\newcommand{\Cjkp}{{C_{j,k}^p}}
\newcommand{\Cjke}{{C_{j,k}^E}}
\newcommand{\Cjkw}{{C_{j,k}^W}}
\newcommand{\Cjks}{{C_{j,k}^S}}
\newcommand{\Cjkn}{{C_{j,k}^N}}

\newcommand{\xjpthreehalf}{{x_{j+\frac{3}{2}}}}

\newcommand{\ujp}{{u_{j+1}}}
\newcommand{\ujpt}{{\tilde{u}_{j+1}}}
\newcommand{\uj}{{u_j}}
\newcommand{\ujph}{{u_{j+\frac{1}{2}}}}
\newcommand{\ujphp}{{u_{j+\frac{1}{2}}^+}}
\newcommand{\ujphm}{{u_{j+\frac{1}{2}}^-}}
\newcommand{\ujmh}{{u_{j-\frac{1}{2}}}}
\newcommand{\ujt}{{\tilde{u}_j}}
\newcommand{\ujm}{{u_{j-1}}}
\newcommand{\cujph}{{(cu)_{j+\frac{1}{2}}}}
\newcommand{\cujmh}{{(cu)_{j-\frac{1}{2}}}}
\newcommand{\cjph}{{c_{j+\frac{1}{2}}}}

\newcommand\eref[1]{(\ref{#1})}

%%%%%%%%%%%%%%%%%%%%%%%%%%%%%%%%%%
%\usepackage{cleveref}
%\usepackage{refcheck}
%%% Infrastructure
\makeatletter
\newcommand{\refcheckize}[1]{%
  \expandafter\let\csname @@\string#1\endcsname#1%
  \expandafter\DeclareRobustCommand\csname relax\string#1\endcsname[1]{%
    \csname @@\string#1\endcsname{##1}\wrtusdrf{##1}}%
  \expandafter\let\expandafter#1\csname relax\string#1\endcsname
}
\makeatother
%%%

%%% Now we add the reference commands we want refcheck to be aware of
%\refcheckize{\eref}

\newcommand*\xbar[1]{%
  \hbox{%
    \vbox{%
      \hrule height 0.5pt % The actual bar
      \kern0.4ex%         % Distance between bar and symbol
      \hbox{%
        \kern-0.05em%      % Shortening on the left side
        \ensuremath{#1}%
        \kern-0.00em%      % Shortening on the right side
      }%
    }%
  }%
}

%%%%%%%%%%%%%%%%%%%%%%%%%%%%%%%%%%%%

%%%%%%%%%%%%%%%%%%%%%%%%%%%%%%%%%%%%
%LORENZO'S COMMANDS


\newcommand{\lop}[0]{\mathcal{L}}
\newcommand{\lopd}[0]{\mathcal{L}_\Delta}
\newcommand{\lopdt}[0]{\mathcal{L}_{\Delta}}
\newcommand{\LL}{\lopdt}
\newcommand{\usol}[0]{\underline{\uvec{q}}_\Delta}
\newcommand{\usoldt}[0]{\underline{\uvec{q}}_\Delta}


\newcommand{\up}[0]{\underline{\uvec{u}}^{(p)}}

\newcommand{\usoldto}[0]{\tilde{\underline{\uvec{q}}}_\Delta}

\newcommand{\uapp}[0]{\uvec{u}_h}
\newcommand{\wapp}[0]{w_h}

\newcommand{\massmatrix}[0]{\mathcal{M}}

\newcommand{\tess}[0]{\mathcal{T}_h}


\newcommand{\uvec}[2][3]{\boldsymbol{#2\mkern-#1mu}\mkern#1mu}


\newcommand{\res}[0]{\textbf{R}}
\newcommand\norm[1]{\left\lVert#1\right\rVert}
\newcommand\abs[1]{\left\lvert#1\right\rvert}

\newcommand{\flux}[0]{\boldsymbol{F}}
\newcommand{\source}[0]{\boldsymbol{S}}
\newcommand{\ST}[0]{\boldsymbol{ST}_i^K}
\newcommand{\extra}[0]{\boldsymbol{ST}_i}


\newcommand{\elres}[0]{\uvec{\Phi}^K(\uapp)}
\newcommand{\noderes}[0]{\uvec{\Phi}^K_i(\uapp)}
\newcommand{\cund}[0]{\underline{\uvec{c}}}
\newcommand{\lopdi}[0]{\mathcal{L}_{\Delta,i}}
\newcommand{\csoldt}[0]{\underline{\uvec{c}}_\Delta}
\newcommand{\basis}[0]{\uvec{v}}
\newcommand{\mis}[0]{\mu}
\def\restriction#1#2{\mathchoice
	{\setbox1\hbox{${\displaystyle #1}_{\scriptstyle #2}$}
		\restrictionaux{#1}{#2}}
	{\setbox1\hbox{${\textstyle #1}_{\scriptstyle #2}$}
		\restrictionaux{#1}{#2}}
	{\setbox1\hbox{${\scriptstyle #1}_{\scriptscriptstyle #2}$}
		\restrictionaux{#1}{#2}}
	{\setbox1\hbox{${\scriptscriptstyle #1}_{\scriptscriptstyle #2}$}
		\restrictionaux{#1}{#2}}}
\def\restrictionaux#1#2{{#1\,\smash{\vrule height .8\ht1 depth .85\dp1}}_{\,#2}} 
\newcommand\swapifbranches[3]{#1{#3}{#2}}
\newcommand{\R}{\mathbb{R}}
\newcommand{\redun}[1]{{\color{red} {\large \textit{Is the following redundant?}} #1}}
\newcommand{\bbu}{\underline{\uvec{u}}}
\newcommand{\bU}{\uvec{U}}
\newcommand{\bbU}{\underline{\uvec{U}}}
\newcommand{\bbr}{\underline{\uvec{r}}}
\newcommand{\by}{\uvec{y}}
\newcommand{\bby}{\underline{\uvec{y}}}
\newcommand{\bG}{{\uvec{G}}}
\newcommand{\bbt}{\underline{t}}
\newcommand{\vecz}{\underline{0}}
\newcommand{\matz}{\underline{\underline{0}}}
\newcommand{\vecbeta}{\underline{\beta}}
\newcommand{\CIP}{\text{CIP}}
\newcommand{\OSS}{\text{OSS}}
\newcommand{\CFL}{\text{CFL}}

\newcommand{\undu}{\underline{\uvec{u}}}
\newcommand{\undy}{\underline{\uvec{y}}}
\newcommand{\undspacetilde}{\widetilde{\underline{\uvec{\phi}}}}
\newcommand{\spacestuff}{\uvec{\phi}}
\newcommand{\undv}{\underline{\uvec{v}}}
\newcommand{\undw}{\underline{\uvec{w}}}
\newcommand{\undr}{\underline{\uvec{r}}}
\newcommand{\undz}{\underline{\uvec{z}}}
\newcommand{\lopdtp}[0]{\mathcal{L}_{\Delta,p}}
\newcommand{\usolp}[0]{\underline{\uvec{q}}_{\Delta}^{(p)}}
\newcommand{\Xp}[0]{X^{(p)}}
\newcommand{\Yp}[0]{Y^{(p)}}
\newcommand{\uex}[0]{\underline{\uvec{u}}_{ex}}
\newcommand{\uexp}[0]{\underline{\uvec{u}}_{ex}^{(p)}}
\newcommand{\lopdLp}[0]{\mathcal{L}_\Delta^{1,(p)}}
\newcommand{\lopdHp}[0]{\mathcal{L}_\Delta^{2,(p)}}
\newcommand{\embep}[0]{\mathcal{E}^{(p)}}
\newcommand{\projp}[0]{\Pi^{(p)}}
\newcommand{\trialspace}[1]{\varphi_{#1}} %phi space _i
\newcommand{\trialspaceiter}[2]{\varphi_{#1}^{#2}} %phi space _i
\newcommand{\trialtime}[1]{\psi^{#1}} %psi time ^m
\newcommand{\trialtimeiter}[2]{\psi^{#1,#2}} %psi time ^m
%THE FOLLOWING IS  ^l TO AVOID CONFUSION WITH u_n
\newcommand{\trialspacetime}[1]{\vartheta^{#1}} %theta space time 
\newcommand{\trialspacetimeiter}[2]{\vartheta^{#1,#2}}

%ADER-FV
\newcommand{\trialspaceFVp}[1]{\varphi_{#1}} %predictor
\newcommand{\trialspaceFVc}[1]{\lambda_{#1}} %corrector 
\newcommand{\APNPM}{{ADER-$\mathbb P_N \mathbb P_M$ }}

%Temperature
\newcommand{\Temp}{T}


\newcommand{\E}{\boldsymbol{E}}

\newcommand{\pnpm}[2]{\mathbb P_{#1}\mathbb P_{#2}}  %PNPM schemes

\newcommand{\lore}[1]{{\leavevmode\color{red} #1}}
\newcommand{\RIIcolor}[1]{{\leavevmode\color{black} #1}}

\newcommand{\hilight}[1]{\colorbox{red}{#1}}

%%%%%%%%%%%%%%%%%%%%%%%%%%%%%%%%%%%%

\newcommand{\varbox}[1]{\makebox[2em][l]{#1}} % adjust 2em to fit your widest var

\newcommand{\rhostar}{\rho^*(t^n)}
\newcommand{\pstar}{p^*(t^n)}

\graphicspath{{Figures/}}

\title{A New Asymptotic-Preserving Dual Formulation Finite-Volume Method for the Compressible Euler Equations}
\author{Alina Chertock\thanks{Department of Mathematics, North Carolina State University, Raleigh, NC 27695, USA;
{\tt chertock@math.ncsu.edu}}, Smadar Karni\thanks{Department of Mathematics, University of Michigan, 48109, USA; {\tt karni@umich.edu}},
Alexander Kurganov\thanks{Department of Mathematica and Shenzhen International Center for Mathematicss, Southern University of Science and
Technology, Shenzhen, 518055, China; {\tt alexander@sustech.edu.cn}}, and Lorenzo Micalizzi\thanks{Department of Mathematics, North Carolina
State University, Raleigh, NC 27695, USA; {\tt lmicali@ncsu.edu}}}
\date{}

\begin{document}
\maketitle

%\vspace*{-0.5cm}
\begin{abstract}
The paper focuses on the development of numerical methods for the compressible Euler equations. It is well-known that if the Mach number is
small, the system becomes stiff and hence explicit schemes suffer from severe time-step restrictions, making them inefficient or even
impractical. Our objective is to develop an asymptotic preserving (AP) scheme that remains uniformly accurate and stable across \lore{all} Mach numbers.
%a wide range of Mach numbers.

Instead of the conservative hyperbolic flux splitting approach, which is widely used to design AP schemes, we \lore{consider a primitive (nonconservative) formulation}
%reformulate the system in the primitive (nonconservative) variables 
and introduce a nonconservative hyperbolic splitting. The resulting %nonconservative
system is
discretized using a semi-implicit \lore{Deferred Correction} approach: the stiff part is handled \lore{(semi-)implicitly} using %the 
second-order central differences, while the
nonstiff part is treated \lore{explicitly} using a second-order path-conservative central-upwind (CU) \lore{discretization}. %scheme. 
A key feature of our method is that
the pressure at each time level is computed by solving a \lore{well-posed} Poisson-type elliptic equation, thereby enforcing the AP property. Simultaneously,
we evolve the conservative form of the system using a semi-discrete CU \lore{scheme}. 
At the end of each \lore{stage of the time discretization},
%step, 
we perform a special
post-processing that selects the appropriate numerical solution depending on the Mach-number. %regime. 
This guarantees that in the
low-Mach-number regime, the solution is obtained by the AP nonconservative scheme, while in higher Mach number regime, \lore{the} sharp and physically 
relevant solution is computed by the conservative CU scheme.

Numerical experiments confirm that the proposed AP scheme achieves the expected second order of accuracy and that the time-step constraint
is independent of the Mach number, making it a robust and efficient alternative to conventional explicit methods.
\end{abstract}

\noindent
{\bf Key words:} Compressible Euler equations; low Mach number; asymptotic preserving (AP) schemes; hyperbolic splitting; semi-implicit
methods; deferred correction.

\medskip
\noindent
{\bf AMS subject classification:} 76M12, 65M08, 65L04, 35B40, 35L60.

\section{Introduction}%\label{sec1}
The paper focuses on \lore{the} compressible Euler equations, which, like any other hyperbolic system of PDEs, are characterized by a finite speed of
propagation. This plays a crucial role in the development of explicit numerical methods, for which a major stability requirement is to keep
the time steps inversely proportional to the maximum wave speed over the entire computational domain.

It is well-known that low-Mach-number flows pose several major challenges for numerical simulations. A distinctive feature of such regimes is
the appearance of both slow material waves, which transport quantities like entropy and vorticity, and fast acoustic waves, whose speeds
scale inversely with the Mach number. As the Mach number decreases, the resulting stiffness imposes severe time-step restrictions on
explicit methods and leads to excessive numerical diffusion, making such schemes inefficient or even impractical for real applications.
Fully-implicit methods can address the stiffness, but have their own drawbacks: they tend to oversmear material waves (see, e.g.,
\cite{boscarino2018all}), require the solution of large nonlinear systems, and may fail to capture the correct solution in the
zero-Mach-number limit.

To overcome these difficulties, a widely adopted strategy is to use either implicit-explicit (IMEX) or semi-implicit (SI) methods based on
conservative hyperbolic flux splitting. This approach decomposes the hyperbolic flux into stiff and nonstiff components in a manner that
preserves the conservative structure of the original system. The fast (stiff) part, associated with acoustic waves, is treated
(semi-)implicitly to relax time-step limitations, while the slow (nonstiff) part is handled explicitly to accurately capture the evolution
of material waves without excessive numerical diffusion.

It is also known that, as the Mach number tends to zero, the compressible Euler equations reduce to the incompressible Euler equations. It is
%therefore 
essential to ensure \lore{that} %IMEX or SI 
numerical schemes also exhibit the same \lore{limit} %ing 
behavior at the discrete level and provide a
consistent discretization of the incompressible Euler equations as the Mach number tends to zero. Schemes that maintain this property are
called asymptotic-preserving (AP). They \lore{were originally introduced to capture} %have their origin in capturing 
steady-state solutions for neutron transport in the diffusive regime
\cite{LMM,LarMor}, but the \lore{specific definition was introduced} %term AP schemes was introduced
in \cite{GJL,jin1999efficient,Kla99} in the context of stiff kinetic equations.
In recent years, AP schemes have been extensively studied and applied for simulating low-Mach-number flows; see, e.g.,
\cite{AllVig,boscarino2019high,boscarino2018all,BDLTV,CGK,cordier2012asymptotic,DJL,degond2011all,DLMV,dimarco2017study,haack2012all,Kle95,arun2014asymptotic,Tang12,TorVaz,zeifang2020novel}
for a non-exhaustive list of references.

All of the aforementioned AP schemes, which were designed for either the isentropic or full compressible Euler system, are based on 
different flux splitting strategies. 
A very simple and robust flux splitting, which was proposed in \cite{haack2012all} for the isentropic 
Euler equations and later extended to the rotating shallow water equations in \cite{LCK}, is \lore{remarkably optimal, in the sense that}
%seem to be rather optimal from the point of view that 
%they very accurately identify and separate 
\lore{it very accurately identifies and separates} a linear stiff pressure term, which is then discretized implicitly. However, extending this 
flux splitting to the full Euler equations presents significant challenges. 

In this paper, we propose \lore{a way to precisely do this, overcoming the aforementioned difficulties}. %n alternative way of accurately identifying and separating a stiff part of the full Euler system: 
\lore{More in detail,} we first
rewrite the %studied 
system \lore{under investigation}
in a nonconservative form and then introduce a nonconservative hyperbolic splitting, which may be naturally
viewed as an extension of the flux splitting from \cite{haack2012all}. We integrate the nonconservative system using a SI method implemented
as follows. The stiff part is approximated semi-implicitly using a simple second-order accurate central-differencing, and the nonstiff part 
is \lore{handled explicitly} %discretized 
using a %the explicit 
second-order path-conservative (PC) central-upwind (CU) \lore{discretization}. %scheme. 
The resulting SI \lore{approach} %discretization 
is then
realized in such a way that the \lore{update of the} pressure % at the new time level 
%is obtained by 
\lore{consists in} solving a Poisson-type elliptic equation: this \lore{is used}
%helps 
to enforce the AP property.

However, the resulting SI method can only be applied in the low-Mach-number regime, where \lore{discontinuities do not occur}. %shock waves are negligible. 
%In the 
\lore{For} large 
%or intermediate 
Mach numbers, solving \lore{nonconservative formulations} %the nonconservative formulation 
of the Euler equations in the presence of discontinuities %will 
\lore{generally leads}
%unavoidably 
%NB: There are ways to solve the nonconservative formulation and still converge to the correct solution see Remi staggered
%lead 
to nonphysical computed solutions, as was demonstrated in \cite{abgrall2010comment,HouLeF}. We therefore apply a dual
formulation (DF) approach, which was recently introduced in \cite{CFKM} (for other recent works on dual formulation
methods, we refer the reader to \cite{Abg2023,ACKM,AbgLiu,pidatella2019semi}), and solve the nonconservative and conservative formulations simultaneously.
The latter one is discretized \lore{in a fully-explicit manner} using the second-order semi-discrete CU discretization from \cite{KLin}. %combined with the explicit realization of the implicit trapezoidal ODE solver. 
This way, at each \lore{stage of the considered Deferred Correction (DeC) time discretization}~\cite{offner2025analysis} %step, 
we will evolve two copies of the computed solution: one is AP, but
nonconservative, while the second one is conservative, but non-AP. Hence, 
%upon the completion of each time step, 
we post-process the obtained solutions to automatically ensure that in the low-Mach-number regime, the overall numerical solution is obtained by the AP
nonconservative scheme, while in the large (intermediate)-Mach-number regime, the solution reduces to the sharp and conservative solution obtained by the CU scheme.

The \lore{rest of the} paper is organized as follows. In \S\ref{sec2}, we give the necessary preliminaries: we introduce the governing equations, namely, non-dimensional conservative and
nonconservative (primitive) formulations of the full Euler equations, discuss their zero-Mach-number limit, and briefly review the
\lore{considered DF framework}. In \S\ref{sec4}, we introduce %a 
\lore{the} novel AP dual formulation-finite volume (DF-FV) scheme for compressible Euler equations, \lore{providing rigorous proof for its AP character}. In \S\ref{sec6}, we
demonstrate the performance of the \lore{proposed} scheme on a number of challenging numerical examples. Finally, concluding remarks can be found
in \S\ref{sec5}.

\section{Preliminaries}\label{sec2}
%The goal of this paper is to develop a new %AP 
%scheme for the compressible Euler equations that remains robust and accurate across %a wide range of 
%\lore{all} Mach numbers and, in particular, yields a consistent approximation of the incompressible limit as the Mach number tends to zero, \lore{i.e., which retains the AP property}. 
%Three key ingredients will be essential in designing this scheme:
\lore{The main goal of this section is to provide the necessary basic ingredients for the presentation of the proposed scheme. Specifically, we will describe:}
\begin{asparaitem}
\item
a \lore{nonconservative} reformulation of the Euler equations in terms of primitive variables, which allows for a natural decomposition \lore{characterized by stiff and nonstiff terms, in \S\ref{sec21}}; % : This will facilitate an efficient SI time integration strategy;
\item
a formal asymptotic analysis of the studied \lore{equations} in the low-Mach-number regime, \lore{providing} the incompressible system that the AP scheme must accurately \lore{capture} in the zero-Mach-number limit, \lore{in \S\ref{sec22}};
%: This will provide the correct incompressible
%limiting system that the AP scheme must accurately approximate in the zero-Mach-number limit;
\item 
\lore{the adopted} DF framework, in which both conservative and nonconservative %(primitive)
formulations of the %studied 
system \lore{under investigation} are numerically solved
simultaneously \lore{exploiting} %utilizing 
the advantages of each of them in the corresponding Mach-number regimes.
\end{asparaitem}

\subsection{Conservative and Primitive Formulations}\label{sec21}
After suitable non-dimensionalization and rescaling, the two-dimensional (2-D) compressible Euler equations can be written in the
conservative form as
\begin{align}
&\rho_t+\bnabla\!\cdot\!(\rho\bm u)=0,\label{2.1a}\\
&(\rho\bm u)_t+\bnabla\!\cdot\!(\rho\bm u\otimes\bm u)+\frac{1}{\ve^2}\nabla p=\bm0,\label{2.1b}\\
&E_t+\bnabla\!\cdot\!((E+p)\bm u)=0.\label{2.1c}
\end{align}
Here, $\rho$, $\bm u:=(u,v)^\top$, and $E$ denote the density, velocity, and total energy, respectively, $p$ is the pressure, $\ve$ is the 
reference Mach number, and the system is closed by the equation of state, which, in the case of an ideal gas, reads as
\begin{equation}
E=\frac{p}{\gamma-1}+\frac{\ve^2}{2}\rho(u^2+v^2),
\label{2.1d}
\end{equation}
with $\gamma$ being the specific heat ratio. This system is hyperbolic and features acoustic waves traveling with \lore{(maximum)} speed $\norm{\bm u}_2+c$, where
$c$ is \lore{the} speed of sound given by $c:=\frac{1}{\ve}\sqrt{\gamma p/\rho}$. Notice that in the low-Mach-number regime, the acoustic waves
travel at a high \lore{(maximum)} speed proportional to $1/\ve$.

For the purpose of deriving \lore{our} AP scheme, % for the compressible Euler equations, 
we also consider an equivalent primitive formulation. %, which our AP method will be based on. 
To this end, we rewrite the system \eref{2.1a}--\eref{2.1d} in terms of the primitive variables
$\rho$, $\bm u$, and~$p$:
\begin{align}
&\rho_t+\bnabla\!\cdot\!(\rho\bm u)=0,\label{2.5a}\\
&\bm u_t+(\bm u\!\cdot\!\bnabla)\bm u+\frac{1}{\ve^2\rho}\nabla p=\bm0,\label{2.5b}\\
&p_t+\bm u\!\cdot\!\nabla p+\gamma p\bnabla\!\cdot\!\bm u=0.\label{2.5c}
\end{align}
We emphasize that this formulation is equivalent to the conservative system \eref{2.1a}--\eref{2.1d} only for smooth solutions, and numerical
approximations of \eref{2.5a}--\eref{2.5c} typically converge to nonphysical solutions \lore{when discontinuities are present}; see 
\cite{abgrall2010comment,HouLeF} for a detailed discussion.

\subsection{\lore{Zero-Mach-Number Limit}}\label{sec22}
It is well-known (see, e.g., a formal analysis in \cite{Ala,MetSch}) that in the zero-Mach-number limit the compressible Euler equations
reduce to the incompressible ones. To illustrate this, we examine the formal behavior of the primitive system \eref{2.5a}--\eref{2.5c} as
$\ve\to0$. We substitute the formal expansions
\begin{align*}
\rho&:=\rho^{(0)}(x,y,t)+\ve\rho^{(1)}(x,y,t)\lore{+\ve^2\rho^{(3)}(x,y,t)}+\dots,\\
\bm u&:=\bm u^{(0)}(x,y,t)+\ve\bm u^{(1)}(x,y,t)\lore{+\ve^2\bm u^{(2)}(x,y,t)}+\dots,\\ 
p&:=p^{(0)}(x,y,t)+\ve p^{(1)}(x,y,t)\lore{+\ve^2 p^{(2)}(x,y,t)}+\dots,
\end{align*}
into \eref{2.5a}--\eref{2.5c} and collect terms by powers of $\ve$. This yields
\begin{equation}
\begin{aligned}
{\cal O}(\ve^{-2}):\quad&\nabla p^{(0)}=\uvec{0},\\
{\cal O}(\ve^{-1}):\quad&\nabla p^{(1)}=\uvec{0},\\
{\cal O}(1):\quad&\rho^{(0)}_t+\bnabla\!\cdot\!(\rho^{(0)}\bm u^{(0)})=0,\\
&\bm u^{(0)}_t+(\bm u^{(0)}\!\cdot\!\bnabla)\bm u^{(0)}+\frac{1}{\rho^{(0)}}\nabla p^{(2)}=\bm0,\\
&p^{(0)}_t+\gamma p^{(0)}\bnabla\!\cdot\!\bm u^{(0)}=0,\\
{\cal O}(\ve):\quad&p^{(1)}_t+\gamma\big(p^{(1)}\bnabla\!\cdot\!\bm u^{(0)}+p^{(0)}\bnabla\!\cdot\!\bm u^{(1)}\big)=0.
\end{aligned}
\label{3.1}
\end{equation}

It follows from the first two equations in \eref{3.1} that $p^{(0)}(x,y,t)=p^{(0)}(t)$ and $p^{(1)}(x,y,t)=p^{(1)}(t)$ are \lore{spatially} uniform.
One can also show that both $p^{(0)}$ and $p^{(1)}$ are independent of time, provided \lore{that} the following Dirichlet-type boundary condition holds:
\begin{equation}
p\equiv\,\xbar p_0+\ve\,\xbar p_1 \lore{+\ve^2 p_2}+\dots ,\quad\forall(x,y)\in\partial\Omega,\quad \lore{\forall t\in [0,T_f]},
\label{3.2}			
\end{equation}
where $\xbar p_0\in \mathbb{R^+}$ and $\xbar p_1\in\mathbb{R}$ are constants, \lore{$p_k:=p_k(x,y,t)$ with $k\geq 2$ are bounded functions}, $\Omega$ is the spatial domain with boundary $\partial\Omega$, \lore{and $T_f$ is the final time}. 
%Moreover, the boundary condition \eref{3.2} 
\lore{In fact, such a boundary condition} implies $p^{(0)}\equiv\,\xbar p_0$ and $p^{(1)}\equiv\,\xbar p_1$. Using this in the $p^{(0)}$-equation in
\eref{3.1}, one concludes that $\bnabla\!\cdot\!\bm u^{(0)}=0$, which further implies $\bnabla\!\cdot\!\bm u^{(1)}=0$ thanks to the $p^{(1)}$-equation.
%
Hence, the zero-Mach-number \lore{limit} equations are
\begin{align}
&\rho^{(0)}_t+\bnabla\!\cdot\!(\rho^{(0)}\bm u^{(0)})=0,\label{3.3}\\
&\bm u^{(0)}_t+(\bm u^{(0)}\!\cdot\!\bnabla)\bm u^{(0)}+\frac{1}{\rho^{(0)}}\nabla p^{(2)}=\bm0,\label{3.4}\\
&p^{(0)}\equiv\,\xbar p_0,\quad p^{(1)}\equiv\,\xbar p_1,\label{3.5_p}\\
&\bnabla\!\cdot\!\bm u^{(0)}=0,\quad \bnabla\!\cdot\!\bm u^{(1)}=0.\label{3.5_u}
\end{align}


\subsection{\lore{Dual Formulation (DF) Framework}}\label{sec4a}
As outlined above, a new AP scheme for the compressible Euler equations will be constructed within the DF framework, in which both
conservative and primitive formulations are numerically solved simultaneously. 
\lore{The conservative form ensures proper handling of
discontinuities, % and conservation properties, 
while the primitive form is used to achieve the AP property in the zero-Mach-number limit.} %stabilize the scheme and manage the terms that become dominant in the stiff case. 
While the DF methodology is not specifically developed to handle multiscale features, it lays the groundwork
for the AP scheme developed in subsequent sections by enabling an efficient and accurate treatment of the studied Euler equations in both
compressible and nearly incompressible flow regimes.

We consider a general 2-D hyperbolic system of conservation laws,
\begin{equation}
\bm U_t+\bm F(\bm U)_x+\bm G(\bm U)_y=\bm0,
\label{4.1f}
\end{equation}
where $\bm U$ is the vector of conservative variables and $\bm F$ and $\bm G$ are fluxes, and rewrite it in an equivalent nonconservative
form
\begin{equation}
\bm V_t+\widetilde{\bm F}(\bm V)_x+\widetilde{\bm G}(\bm V)_y=B(\bm V)\bm V_x+C(\bm V)\bm V_y,
\label{4.2f}
\end{equation}
where $\bm V$ is the vector of nonconservative variables, $\widetilde{\bm F}$ and $\widetilde{\bm G}$ are the corresponding fluxes, and
$B(\bm V)\bm V_x$ and $C(\bm V)\bm V_y$ are the nonconservative product terms.

\lore{The key idea of the DF approach is to construct semi-discretizations of \eref{4.1f} and \eref{4.2f}, which are evolved in time simultaneously.} A crucial step in the DF methods is \lore{the} post-processing, in which the evolved values of the
nonconservative solution $\bm V$ \lore{are} replaced with a more reliable approximation after the update. This step is necessary
because long-term \lore{evolutions} of $\bm V$ by \lore{directly} solving the nonconservative system \eref{4.2f} may lead to \lore{nonphysical solutions in presence of discontinuities, }
%discontinuities in the solution, 
which typically appear when the studied Euler system is considered in the compressible (large/medium-Mach-number) regime.

\lore{The} post-processing can be described as follows. After computing the solutions of \eref{4.1f} and \eref{4.2f} \lore{corresponding to a certain time level $t^*$}, % from the time level $t$ to the next time level $t+\dt$ 
the evolved values of \lore{$\bm V(t^*)$} are replaced with
\begin{equation}
\lore{r\big(\bm V(\bm U(t^*)),\bm V(t^*)\big),}
\label{4.5f}
\end{equation}
where $r$ is a suitable replacement function and $\bm V(\bm U)$ is a conservative-primitive variable transformation. In the simplest
nonstiff case, considered in \cite{CFKM}, one sets
\begin{equation}
\lore{r\big(\bm V(\bm U(t^*)),\bm V(t^*)\big)=\bm V(\bm U(t^*)).}
\label{4.6f}
\end{equation}
However, in the development of the AP scheme below, we will modify the post-processing \eref{4.6f} by taking an appropriate function $r$ in
\eref{4.5f} to ensure that in the nearly incompressible (low-Mach-number) regime the AP $\bm V$-solution is not overwritten by the non-AP
conservative one.

\lore{A crucial difference with respect to~\cite{CFKM} is that here we perform the post-processing at each time stage of the considered time discretization, while, in the aforementioned reference, this is done upon completion of each full time step. Another difference is that here we consider a finite-volume (FV) framework, hence regarding our scheme as AP DF-FV, while, in~\cite{CFKM} a finite-difference framework has been considered. The strategy is however general and can be applied to many other numerical approaches.}

\section{Novel AP Scheme for Compressible Euler Equations}\label{sec4}
Building on the \lore{DF} framework described in \S\ref{sec4a}, we now present a novel AP scheme designed for the compressible Euler equations
across %a wide range of 
\lore{all} Mach number regimes, from fully compressible to nearly incompressible flows. The proposed method couples the
conservative \eref{2.1a}--\eref{2.1d} and primitive \eref{2.5a}--\eref{2.5c} formulations of the system, ensuring stability, accuracy, and
consistency with the %correct 
\lore{analytical} asymptotic behavior as $\ve\to0$.

In this section, we provide a complete description of the proposed space-time discretization, starting with the primitive system
\eref{2.5a}–\eref{2.5c}. In \S\ref{sec411}, we outline its temporal integration, which is based on a new hyperbolic splitting and an SI
approach. 
In \S\ref{sec32}, we present a \lore{fully-discrete} second-order AP scheme for the primitive system, \lore{while}, %and 
in \S\ref{sec42}, we describe
the semi-discrete CU scheme \lore{employed} for the conservative system. \S\ref{sec422} is devoted to the clarification 
of important implementation details.
Finally, in \S\ref{sec43}, we present \lore{the Mach-number dependent} post-processing strategy used to reconcile primitive and conservative variables.

\subsection{Novel AP Time Discretization of the Primitive System}\label{sec411}
We begin by providing a precise definition of %what it means for a time discretization to be 
\lore{asymptotic-preserving (AP) time discretization in the context of
the low-Mach-number limit.}
\begin{definition}[AP time discretization]\label{def41}
	\lore{A time discretization is AP if, assuming}
	\begin{itemize}
		\item \lore{an expansion for the solution at time $t^n$ of the type}
		\begin{equation}\label{eq:expansion_tn}
			\begin{aligned}
				\rho^n &:= \rho^{(0),n} + \varepsilon \rho^{(1),n}
					+ \varepsilon^2 \rho^{(2),n} + \dots, \\
				\bm u^n &:= \bm u^{(0),n} + \varepsilon \bm u^{(1),n}
					+ \varepsilon^2 \bm u^{(2),n} + \dots, \\
				p^n &:= p^{(0),n} + \varepsilon p^{(1),n}
					+ \varepsilon^2 p^{(2),n} + \dots,
			\end{aligned}
		\end{equation}
		\lore{compatible with the asymptotic limit, i.e., satisfying}
	%Assume that the computed solution can be expanded as
		\begin{align}
			&p^{(0),n}\equiv\,\xbar p_0,\quad p^{(1),n}\equiv\,\xbar p_1,\label{eq:p_const}\\
			&\bnabla\!\cdot\!\bm u^{(0),n}=0,\quad \bnabla\!\cdot\!\bm u^{(1),n}=0,\label{eq:u_div_free}
		\end{align}
		\item 	\lore{the Dirichlet-type boundary condition~\eref{3.2}},
		\end{itemize}	
%with $p^{(0),n}\equiv\,\xbar p_0>0$, $p^{(1),n}\equiv\,\xbar p_1>0$, and both $\bm u^{(0),n}$ and $\bm u^{(1),n}$ being divergence-free,
%that is,
%\begin{equation}
%\bnabla\!\cdot\!\bm u^{(0),n}=0,\quad\bnabla\!\cdot\!\bm u^{(1),n}=0.
%\label{4.2}
%\end{equation}
%In addition, assume that the boundary condition \eref{3.2} is satisfied, namely,
%$$
%p^n\equiv\,\xbar p_0+\ve\,\xbar p_1+{\cal O}(\ve^2),\quad\forall(x,y)\in\partial\Omega.
%$$
		\lore{then, it produces an approximation $(\rho^{n+1},\bm u^{n+1},p^{n+1})$ at time $t^{n+1}$ consistent with the asymptotic limit, namely, if the obtained solution admits an expansion of the type~\eqref{eq:expansion_tn} satisfying conditions~\eqref{eq:p_const} and~\eqref{eq:u_div_free}, and if such data are associated with a consistent discretization of~\eqref{3.3} and~\eqref{3.4}.}
%Let us consider a one-step time discretization that allows to compute the solution at time $t^{n+1}=t^n+\dt$. We say that such time
%discretization is AP if it yields a consistent discretization of \eref{3.3}--\eref{3.5} as $\ve\to0$ and the solution
%$(\rho^{n+1},\bm u^{n+1},p^{n+1})$ at $t^{n+1}$ admits the same asymptotic expantion as in \eref{4.1} with $p^{(0),n+1}\equiv\,\xbar p_0$,
%$p^{(1),n+1}\equiv\,\xbar p_1$, and both $\bm u^{(0),n+1}$ and $\bm u^{(1),n+1}$ being divergence-free.
\end{definition}

To construct an AP time discretization, we first perform a hyperbolic splitting of the primitive system \eref{2.5a}–\eref{2.5c}, separating
the stiff pressure-driven terms from the nonstiff convective terms. This splitting enables the use of an SI integration strategy in which
the stiff and nonstiff terms are treated semi-implicitly and explicitly, respectively, ensuring uniform stability and asymptotic consistency
as $\ve\to0$.

\subsubsection{A New Hyperbolic Splitting}
We follow the idea from \cite{haack2012all,LCK} and split the nonconservative system into two parts corresponding to the slow and fast
dynamics as follows. We first define the time-dependent variables 
\begin{equation}
\rho_{\max}(t):=\max\limits_{(x,y)\in\Omega}\rho(x,y,t),\quad p_{\min}(t):=\min\limits_{(x,y)\in\Omega}p(x,y,t),
\label{4.3}
\end{equation}
and then add and subtract $\frac{1}{\ve^2\rho_{\max}}\nabla p$ and $\gamma p_{\min}\bnabla\cdot\bm u$ from \eref{2.5b} and \eref{2.5c},
respectively, to rewrite system \eref{2.5a}--\eref{2.5c} as follows:
\begin{align}
&\rho_t+\bnabla\!\cdot\!(\rho\bm u)=0,\label{2.1}\\
&\bm u_t+(\bm u\!\cdot\!\bnabla)\bm u+\frac{\rho_{\max}-\rho}{\ve^2\rho\rho_{\max}}\,\nabla p=-\frac{1}{\ve^2\rho_{\max}}
\nabla p,\label{2.2}\\[0.5ex]
&p_t+\bm u\!\cdot\!\nabla p+\gamma(p-p_{\min})\bnabla\!\cdot\!\bm u=-\gamma p_{\min}\bnabla\!\cdot\!\bm u.\label{2.3}
\end{align}
This system can be put in the following vector form:
\begin{equation}
\bm V_t+\widetilde{\mF}(\bm V)_x+\widetilde{\mG}(\bm V)_y=\widetilde B(\bm V)\bm V_x+\widetilde C(\bm V)\bm V_y+\widehat B(\bm V)\bm V_x+
\widehat C(\bm V)\bm V_y,
\label{4.47}
\end{equation}
where $\bm V:=(\rho,u,v,p)^\top$, the nonlinear nonstiff (slow dynamics) part consists of the fluxes
\begin{equation*}
\widetilde{\mF}(\bm V)=\Big(\rho u,\frac{u^2}{2},0,0\Big)^\top\quad\mbox{and}\quad
\widetilde{\mG}(\bm V)=\Big(\rho v,0,\frac{v^2}{2},0\Big)^\top,
\end{equation*}
and \lore{of} the nonstiff nonconservative terms % are 
$\widetilde B(\bm V)\bm V_x+\widetilde C(\bm V)\bm V_y$ with matrices %the
\begin{equation*}
\widetilde B=-\begin{pmatrix}0&0&0&0\\0&0&0&\dfrac{\rho_{\max}-\rho}{\ve^2\rho\rho_{\max}}\\0&v&0&0\\0&\gamma(p-p_{\min})&0&u\end{pmatrix}
\quad\mbox{and}\quad
\widetilde C=-\begin{pmatrix}0&0&0&0\\0&0&u&0\\0&0&0&\dfrac{\rho_{\max}-\rho}{\ve^2\rho\rho_{\max}}\\0&0&\gamma(p-p_{\min})&v\end{pmatrix},
\end{equation*}
while, the linear stiff (fast dynamics) part consists of the stiff nonconservative terms $\widehat B(\bm V)\bm V_x+\widehat C(\bm V)\bm V_y$
with matrices %the
\begin{equation*}
\widehat B=-\begin{pmatrix}0&0&0&0\\0&0&0&\dfrac{1}{\ve^2\rho_{\max}}\\0&0&0&0\\0&\gamma p_{\min}&0&0\end{pmatrix}
\quad\mbox{and}\quad
\widehat C=-\begin{pmatrix}0&0&0&0\\0&0&0&0\\0&0&0&\dfrac{1}{\ve^2\rho_{\max}}\\0&0&\gamma p_{\min}&0\end{pmatrix}.
\end{equation*}

We stress that the subsystem
$\bm V_t+\widetilde{\mF}(\bm V)_x+\widetilde{\mG}(\bm V)_y=\widetilde B(\bm V)\bm V_x+\widetilde C(\bm V)\bm V_y$ is indeed nonstiff as the
eigenvalues of the matrices
$\frac{\partial\widetilde{\mF}}{\partial\bm V}(\uvec{V})-\widetilde B(\bm V)$ and
$\frac{\partial\widetilde{\mG}}{\partial\bm V}(\uvec{V})-\widetilde C(\bm V)$, are $\left\{u\pm\tilde c,u,u\right\}$ and
$\left\{v\pm\tilde c,v,v\right\}$, respectively, with 
\begin{equation}
\tilde c:=\frac{1}{\ve}\sqrt{\gamma\frac{(\rho_{\max}-\rho)(p-p_{\min})}{\rho\rho_{\max}}},
\label{4.48f}
\end{equation}
which are real and of size ${\cal O}(1)$ thanks to the definitions of $\rho_{\max}$ and $p_{\min}$ in \eref{4.3} and \lore{to} the asymptotic
analysis in \S\ref{sec22}, which ensure that
\begin{equation}
0\le\rho_{\max}-\rho={\cal O}(1)\quad\mbox{and}\quad0\le p-p_{\min}={\cal O}(\ve^2).
\label{4.7}
\end{equation}

In the next subsection, we will utilize this splitting and design an AP time discretization %by 
\lore{based on} an explicit approximation of the nonstiff 
subsystem and \lore{an} SI discretization of the stiff terms on the right-hand sides (RHSs) of \eref{2.2} and \eref{2.3}.

\subsubsection{First-Order AP SI Time Discretization}
The simplest first-order AP SI time discretization of the system \eref{4.47} reads as
\begin{equation}
\begin{aligned}
&\frac{\rho^{n+1}-\rho^n}{\dt}+\bnabla\!\cdot\!(\rho^n\bm u^n)=0,\\
&\frac{\bm u^{n+1}-\bm u^n}{\dt}+(\bm u^n\!\cdot\!\bnabla)\bm u^n+\frac{\rho_{\max}^n-\rho^n}{\ve^2\rho^n\rho_{\max}^n}\,\nabla p^n+
\frac{1}{\ve^2\rho_{\max}^n}\nabla p^{n+1}=\bm0,\\
&\frac{p^{n+1}-p^n}{\dt}+\bm u^n\!\cdot\!\nabla p^n+\gamma(p^n-p_{\min}^n)\bnabla\!\cdot\!\bm u^n+\gamma p_{\min}^n
\bnabla\!\cdot\!\bm u^{n+1}=0,
\end{aligned}
\label{4.8}
\end{equation}
which can also be written in the following vector form:
\begin{equation*}
\frac{\bm V^{n+1}-\bm V^n}{\dt}+\bm{{\cal R}}^n+\bm{{\cal L}}^{n,n+1}=\bm0,
\end{equation*}
where $\bm V^n:=(\rho^n,\bm u^n,p^n)^\top\approx\bm V(t^n)$, $\rho_{\max}^n:=\rho_{\max}(t^n)$, $p_{\min}^n:=p_{\min}(t^n)$, and 
\begin{align}
&\begin{aligned}
\bm{{\cal R}}^n&:=\widetilde{\mF}(\bm V^n)_x+\widetilde{\mG}(\bm V^n)_y-\widetilde B(\bm V^n)\bm V_x^n-\widetilde C(\bm V^n)\bm V_y^n\\
&\hspace*{0.15cm}=\begin{pmatrix}({\cal R}^\rho)^n\\(\bm{{\cal R}}^{\bm u})^n\\({\cal R}^p)^n\end{pmatrix}=
\begin{pmatrix}\lore{\bnabla}\!\cdot\!\rho^n\bm u^n\\(\bm u^n\!\cdot\!\bnabla)\bm u^n+\dfrac{\rho_{\max}^n-\rho^n}{\ve^2\rho^n\rho_{\max}^n}\lore{\nabla p^n}\\[1.5ex]
\bm u^n\!\cdot\!\nabla p^n+\gamma(p^n-p_{\min}^n)\bnabla\!\cdot\!\bm u^n\end{pmatrix},
\end{aligned}\label{3.12f}\\
&\bm{{\cal L}}^{n,n+1}:=-\widehat B(\bm V^n)\bm V_x^{n+1}-\widehat C(\bm V^n)\bm V_y^{n+1}=
\begin{pmatrix}({\cal L}^\rho)^{n,n+1}\\(\bm{{\cal L}}^{\bm u})^{n,n+1}\\({\cal L}^p)^{n,n+1}\end{pmatrix}
=\begin{pmatrix}0\\\dfrac{\nabla p^{n+1}}{\ve^2\rho_{\max}^n}\\[1.5ex]
\gamma p_{\min}^n\bnabla\!\cdot\!\bm u^{n+1}\end{pmatrix}.\label{3.13f}
\end{align}

Notice that in \eref{3.13f}, $\bm{{\cal L}}^{n,n+1}$ is discretized in an SI (rather than \lore{fully-}implicit) manner, where both $\widehat B$ and
$\widehat C$ are evaluated at $\bm V^n$ (and not at $\bm V^{n+1}$), which 
%helps to avoid the necessity of 
\lore{prevents from} numerically solving complicated
systems of nonlinear algebraic equations.

We shall now prove that the time discretization \eref{4.8} is indeed AP, provided that the time step is computed based on the wave speeds of
the nonstiff subsystem, that is, provided that
\begin{equation}
\dt=K_{\rm CFL}\min\left\{\frac{\dx}{\max\limits_{(x,y)\in\Omega}\big(|u|+\tilde c\big)},\,
\frac{\dy}{\max\limits_{(x,y)\in\Omega}\big(|v|+\tilde c\big)}\right\},
\label{4.10}
\end{equation}
where $K_{\rm CFL}$ is a CFL number and $\dx$ and $\dy$ are mesh sizes used in the spatial discretization. Notice that selecting the time
step $\dt$ according to \eref{4.10} makes it \lore{asymptotically} independent of $\ve$ as, according to \eref{4.48f}--\eref{4.7}, $\tilde c={\cal O}(1)$.

\begin{theorem}\label{thm31}
The first-order SI time discretization \eref{4.8} is AP according to Definition \ref{def41}, provided that $\dt$ is computed \lore{as in} %according to
\eref{4.10}.
\end{theorem}
\begin{proof}
We begin by formally showing that the computed solution $\rho^{n+1}$, $\bm u^{n+1}$, $p^{n+1}$ admits an expansion of the type \eqref{eq:expansion_tn} \lore{with $(0)$- and $(1)$-pressure and velocity modes satisfying~\eqref{eq:p_const} and~\eqref{eq:u_div_free} respectively}
%with $p^{(0),n}\equiv\,\xbar p_0$, $p^{(1),n}\equiv\,\xbar p_1$, and $\bnabla\!\cdot\!\bm u^{(0),n+1}=\bnabla\!\cdot\!\bm u^{(1),n+1}=0$ 
in the limit as $\ve\to0$. 
We substitute the corresponding expansion of the numerical solution at time $t^n$ into the scheme \eref{4.8} and use \lore{assumptions~\eqref{eq:p_const} and~\eqref{eq:u_div_free}} to obtain
\begin{align}
	&\begin{aligned}
	\rho^{n+1}&=\rho^{(0),n}+\ve \rho^{(1),n}+\ve^2 \rho^{(2),n}-\dt\bnabla\!\cdot\!\big(\rho^{(0),n}\bm u^{(0),n}\big)\\
	&-\ve \dt \big[\bnabla\!\cdot\!\big(\rho^{(1),n}\bm u^{(0),n}\big)+\bnabla\!\cdot\!\big(\rho^{(0),n}\bm u^{(1),n}\big)\big]\\
	&-\ve^2 \dt \big[\bnabla\!\cdot\!\big(\rho^{(2),n}\bm u^{(0),n}\big)+\bnabla\!\cdot\!\big(\rho^{(1),n}\bm u^{(1),n}\big)+\bnabla\!\cdot\!\big(\rho^{(0),n}\bm u^{(2),n}\big)\big]+{\cal O}(\ve^3),\label{4.11}
	\end{aligned}\\
	&\begin{aligned}
\bm u^{n+1}&=\bm u^{(0),n}+\ve \bm u^{(1),n}+\ve^2 \bm u^{(2),n}-\dt \Big[(\bm u^{(0),n}\!\cdot\!\bnabla)\bm u^{(0),n}+\frac{\rho_{\max}^n-\rho^{(0),n}}{\rho^{(0),n}\rho_{\max}^n}\,\nabla p^{(2),n} \Big]
\\&-\ve\dt \Big[(\bm u^{(1),n}\!\cdot\!\bnabla)\bm u^{(0),n}+(\bm u^{(0),n}\!\cdot\!\bnabla)\bm u^{(1),n}-\frac{\rho^{(1),n}}{\rho^{(0),n}\rho_{\max}^n}\,\nabla p^{(2),n}-\frac{\rho^{(0),n}}{\rho^{(0),n}\rho_{\max}^n}\,\nabla p^{(3),n}\Big]\\
&-\ve^2\dt \Big[
(\bm u^{(2),n}\!\cdot\!\bnabla)\bm u^{(0),n}
+(\bm u^{(1),n}\!\cdot\!\bnabla)\bm u^{(1),n}
+(\bm u^{(0),n}\!\cdot\!\bnabla)\bm u^{(2),n} \\
&\qquad
-\frac{\rho^{(2),n}}{\rho^{(0),n}\rho_{\max}^n}\,\nabla p^{(2),n}
-\frac{\rho^{(1),n}}{\rho^{(0),n}\rho_{\max}^n}\,\nabla p^{(3),n}
-\frac{\rho^{(0),n}}{\rho^{(0),n}\rho_{\max}^n}\,\nabla p^{(4),n}
\Big]\\
&-\frac{\dt}{\ve^2 \rho_{\max}^n}\,\nabla p^{n+1}+{\cal O}(\ve^3),\label{4.12}
\end{aligned}\\
p^{n+1}&=\,\xbar p_0+\ve\,\xbar p_1+\ve^2 p^{(2),n}-\ve^2 \dt \bm u^{(0),n}\!\cdot\!\nabla p^{(2),n} -\dt\gamma p_{\min}^n\bnabla\!\cdot\!\bm u^{n+1}+{\cal O}(\ve^3).\label{4.13}
\end{align}
Thanks to the explicit nature of the density update in \eref{4.11}, we conclude that indeed $\rho^{n+1}$ admits the required asymptotic expansion $\rho^{n+1}=\rho^{(0),n+1}+\ve\rho^{(1),n+1}+\ve\rho^{(2),n+1}+\dots$, where the different terms of the expansion are obtained by collecting corresponding powers of $\ve$; for example, ones has
\begin{align*}
	\rho^{(0),n+1}&=\rho^{(0),n}-\dt\,\bnabla\!\cdot\!(\rho^{(0),n}\bm u^{(0),n}),\\
	\rho^{(1),n+1}&=\rho^{(1),n}-\dt \big[\bnabla\!\cdot\!\big(\rho^{(1),n}\bm u^{(0),n}\big)+\bnabla\!\cdot\!\big(\rho^{(0),n}\bm u^{(1),n}\big)\big],\\
	\rho^{(2),n+1}&=\rho^{(2),n}-\dt \big[\bnabla\!\cdot\!\big(\rho^{(2),n}\bm u^{(0),n}\big)+\bnabla\!\cdot\!\big(\rho^{(1),n}\bm u^{(1),n}\big)+\bnabla\!\cdot\!\big(\rho^{(0),n}\bm u^{(2),n}\big)\big].
\end{align*}
% with
%which shows that the density equation in the scheme \eref{4.8} is consistent with \eref{3.3} as $\ve\to0$ since according to \eref{4.10},
%$\dt$ is independent of $\ve$.
In order to show that also $p^{n+1}$ admits an expansion of the same type, we will show that it satisfies a specific well-posed elliptic problem. More in detail, we take the divergence of the velocity equation \eref{4.12} and we substitute
$\bnabla\!\cdot\!\muu^{n+1}$ into the pressure equation \eref{4.13}. Recalling the divergence-free assumption~\eqref{eq:u_div_free}, we obtain
%, take into account , and 
%This yields

$$
\begin{aligned}
p^{n+1}&=\,\xbar p_0+\ve\,\xbar p_1+\ve^2 p^{(2),n}-\ve^2 \dt \bm u^{(0),n}\!\cdot\!\nabla p^{(2),n}\\
&-\dt\gamma p_{\min}^n\bnabla\!\cdot\!\bm u^{(0),n}-\ve\dt\gamma p_{\min}^n\bnabla\!\cdot\!\bm u^{(1),n}-\ve^2\dt\gamma p_{\min}^n\bnabla\!\cdot\!\bm u^{(2),n}\\
&+(\dt)^2\gamma p_{\min}^n\bnabla\!\cdot\!\ \Big[(\bm u^{(0),n}\!\cdot\!\bnabla)\bm u^{(0),n}+\frac{\rho_{\max}^n-\rho^{(0),n}}{\rho^{(0),n}\rho_{\max}^n}\,\nabla p^{(2),n} \Big]\\
&+\ve (\dt)^2\gamma p_{\min}^n\bnabla\!\cdot\!\ \Big[(\bm u^{(1),n}\!\cdot\!\bnabla)\bm u^{(0),n}+(\bm u^{(0),n}\!\cdot\!\bnabla)\bm u^{(1),n}-\frac{\rho^{(1),n}}{\rho^{(0),n}\rho_{\max}^n}\,\nabla p^{(2),n}-\frac{\rho^{(0),n}}{\rho^{(0),n}\rho_{\max}^n}\,\nabla p^{(3),n}\Big]\\
&+\ve^2 (\dt)^2\gamma p_{\min}^n\bnabla\!\cdot\!\ \Big[
(\bm u^{(2),n}\!\cdot\!\bnabla)\bm u^{(0),n}
+(\bm u^{(1),n}\!\cdot\!\bnabla)\bm u^{(1),n}
+(\bm u^{(0),n}\!\cdot\!\bnabla)\bm u^{(2),n} \\
&\qquad
-\frac{\rho^{(2),n}}{\rho^{(0),n}\rho_{\max}^n}\,\nabla p^{(2),n}
-\frac{\rho^{(1),n}}{\rho^{(0),n}\rho_{\max}^n}\,\nabla p^{(3),n}
-\frac{\rho^{(0),n}}{\rho^{(0),n}\rho_{\max}^n}\,\nabla p^{(4),n}
\Big]\\
&+\frac{(\dt)^2\gamma p_{\min}^n}{\ve^2 \rho_{\max}^n}\,\Delta p^{n+1}+{\cal O}(\ve^3),
\end{aligned}
$$
which implies that $p^{n+1}$ is the solution of the following elliptic problem
%$$
%\begin{aligned}
%p^{n+1}&-\frac{(\dt)^2\gamma p_{\min}^n}{\ve^2\rho_{\max}^n}\,\Delta p^{n+1}=
%\,\xbar p_0+\ve\,\xbar p_1+(\dt)^2\gamma p_{\min}^n\bnabla\!\cdot\!\bigg[(\bm u^{(0),n}\!\cdot\!\bnabla)\bm u^{(0),n}\\
%&+\ve\Big((\bm u^{(0),n}\!\cdot\!\bnabla)\bm u^{(1),n}+(\bm u^{(1),n}\!\cdot\!\bnabla)\bm u^{(0),n}\Big)
%+\frac{\rho_{\max}^n-\rho^{(0),n}}{\rho^{(0),n}\rho_{\max}^n}\nabla p^{(2),n}\bigg]+{\cal O}(\ve^2),
%\end{aligned}
%$$
%so that
$$
-\Delta p^{n+1}+\frac{\ve^2\rho_{\max}^n}{(\dt)^2\gamma p_{\min}^n}\,p^{n+1}={\cal O}(\ve^2).
$$
Recalling the assumed boundary condition~\eqref{3.2}, standard theory of perturbed linear operators~\cite{kato1966perturbation} implies that 
\begin{align}
	&p^{n+1}= p^{(0),n+1}+\ve p^{(1),n+1}+\ve^2 p^{(2),n+1}+\dots,\\
    &p^{(0),n+1}=p^{(0),n}\equiv\,\xbar p_0,~~p^{(1),n+1}=p^{(1),n}\equiv\,\xbar p_1.
	\label{4.14}
\end{align}
%which implies $\Delta p^{n+1}={\cal O}(\ve^2)$. This together with the boundary conditions \eref{3.2} results in

At this point, we can easily show that also the updated velocity admits the desired expansion by plugging in the related update the obtained pressure expansion, yielding
$$
\begin{aligned}
\bm u^{n+1}&=\bm u^{(0),n}+\ve \bm u^{(1),n}+\ve^2 \bm u^{(2),n}-\dt \Big[(\bm u^{(0),n}\!\cdot\!\bnabla)\bm u^{(0),n}+\frac{\rho_{\max}^n-\rho^{(0),n}}{\rho^{(0),n}\rho_{\max}^n}\,\nabla p^{(2),n} \Big]
\\&-\ve\dt \Big[(\bm u^{(1),n}\!\cdot\!\bnabla)\bm u^{(0),n}+(\bm u^{(0),n}\!\cdot\!\bnabla)\bm u^{(1),n}-\frac{\rho^{(1),n}}{\rho^{(0),n}\rho_{\max}^n}\,\nabla p^{(2),n}-\frac{\rho^{(0),n}}{\rho^{(0),n}\rho_{\max}^n}\,\nabla p^{(3),n}\Big]\\
&-\ve^2\dt \Big[
(\bm u^{(2),n}\!\cdot\!\bnabla)\bm u^{(0),n}
+(\bm u^{(1),n}\!\cdot\!\bnabla)\bm u^{(1),n}
+(\bm u^{(0),n}\!\cdot\!\bnabla)\bm u^{(2),n} \\
&\qquad
-\frac{\rho^{(2),n}}{\rho^{(0),n}\rho_{\max}^n}\,\nabla p^{(2),n}
-\frac{\rho^{(1),n}}{\rho^{(0),n}\rho_{\max}^n}\,\nabla p^{(3),n}
-\frac{\rho^{(0),n}}{\rho^{(0),n}\rho_{\max}^n}\,\nabla p^{(4),n}
\Big]\\
&-\frac{\dt}{\rho_{\max}^n}\,\nabla p^{(2),n+1}-\ve \frac{\dt}{\rho_{\max}^n}\,\nabla p^{(3),n+1}-\ve^2 \frac{\dt}{\rho_{\max}^n}\,\nabla p^{(4),n+1}+{\cal O}(\ve^3).
\end{aligned}
$$
From this, we indeed get $\bm u^{n+1}= \bm u^{(0),n+1}+\ve \bm u^{(1),n+1}+\ve^2 \bm u^{(2),n+1}+\dots$, with expansion coefficients $\bm u^{(k),n+1}$ obtained collecting the terms in powers of $\ve$ as done for the density; for example
\begin{align*}
	\bm u^{(0),n+1}&=\bm u^{(0),n}-\dt \Big[(\bm u^{(0),n}\!\cdot\!\bnabla)\bm u^{(0),n}+\frac{\rho_{\max}^n-\rho^{(0),n}}{\rho^{(0),n}\rho_{\max}^n}\,\nabla p^{(2),n} \Big]-\frac{\dt}{\rho_{\max}^n}\,\nabla p^{(2),n+1},\\
	\bm u^{(1),n+1}&=\bm u^{(1),n}-\dt \Big[(\bm u^{(1),n}\!\cdot\!\bnabla)\bm u^{(0),n}+(\bm u^{(0),n}\!\cdot\!\bnabla)\bm u^{(1),n}-\frac{\rho^{(1),n}}{\rho^{(0),n}\rho_{\max}^n}\,\nabla p^{(2),n}-\frac{\rho^{(0),n}}{\rho^{(0),n}\rho_{\max}^n}\,\nabla p^{(3),n}\Big]\\
	&-\frac{\dt}{\rho_{\max}^n}\,\nabla p^{(3),n+1},\\
	\bm u^{(2),n+1}&=\bm u^{(2),n}-\dt \Big[
	(\bm u^{(2),n}\!\cdot\!\bnabla)\bm u^{(0),n}
	+(\bm u^{(1),n}\!\cdot\!\bnabla)\bm u^{(1),n}
	+(\bm u^{(0),n}\!\cdot\!\bnabla)\bm u^{(2),n} \\
	&\qquad
	-\frac{\rho^{(2),n}}{\rho^{(0),n}\rho_{\max}^n}\,\nabla p^{(2),n}
	-\frac{\rho^{(1),n}}{\rho^{(0),n}\rho_{\max}^n}\,\nabla p^{(3),n}
	-\frac{\rho^{(0),n}}{\rho^{(0),n}\rho_{\max}^n}\,\nabla p^{(4),n}
	\Big]-\frac{\dt}{\rho_{\max}^n}\,\nabla p^{(4),n+1}.
\end{align*}


%Next, we substitute \eref{4.14} into the the velocity equation \eref{4.12} to obtain
%$$
%\begin{aligned}
%\bm u^{n+1}=\bm u^{(0),n}&-\dt\bigg[(\bm u^{(0),n}\!\cdot\!\bnabla)\bm u^{(0),n}+
%\frac{\rho_{\max}^n-\rho^{(0),n}}{\rho^{(0),n}\rho_{\max}^n}\,\nabla p^{(2),n}+\frac{1}{\rho_{\max}^n}\nabla p^{(2),n+1}\bigg]\\[0.5ex]
%&+\ve\Big[\bm u^{(1),n}-\dt\left((\bm u^{(0),n}\!\cdot\!\bnabla)\bm u^{(1),n}+(\bm u^{(1),n}\!\cdot\!\bnabla)\bm u^{(0),n}\right)\Big]+{\cal O}(\ve^2),
%\end{aligned}
%$$
%which implies that $\bm u^{n+1}=\bm u^{(0),n+1}+\ve\bm u^{(1),n+1}+\dots$ with
%\begin{equation}
%\bm u^{(0),n+1}=\bm u^{(0),n}-\dt\bigg[(\bm u^{(0),n}\!\cdot\!\bnabla)\bm u^{(0),n}+
%\frac{\rho_{\max}^n-\rho^{(0),n}}{\rho^{(0),n}\rho_{\max}^n}\nabla p^{(2),n}\bigg]-\frac{\dt}{\rho_{\max}^n}\nabla p^{(2),n+1}
%\label{4.15}
%\end{equation}
%and
%\begin{equation}
%\bm u^{(1),n+1}=\bm u^{(1),n}-\dt\left((\bm u^{(0),n}\!\cdot\!\bnabla)\bm u^{(1),n}+(\bm u^{(1),n}\!\cdot\!\bnabla)\bm u^{(0),n}\right).
%\label{3.19f}
%\end{equation}
%It should be observed that taking the divergence of \eref{4.15} and \eref{3.19f} implies the divergence-free property at time level
%$t=t^{n+1}$, namely, $\bnabla\!\cdot\!\bm u^{(0),n+1}=\bnabla\!\cdot\!\bm u^{(1),n+1}=0$.

We now need to show that conditions~\eqref{eq:p_const} and~\eqref{eq:u_div_free} hold for the updated solution, along with consistency with~\eqref{3.3} and~\eqref{3.4}.
%
We have already shown that~\eqref{eq:p_const} holds true, see~\eqref{4.14}.
%
Instead, the divergence--free conditions~\eqref{eq:u_div_free} can be deduced from the pressure update~\eqref{4.13}, which in view of the obtained results yields
$$\begin{aligned}
	\ve^2 p^{(2),n+1}=\ve^2 p^{(2),n}&-\ve^2 \dt \bm u^{(0),n}\!\cdot\!\nabla p^{(2),n} -\dt\gamma p_{\min}^n\bnabla\!\cdot\!\bm u^{(0),n+1}\\
	&-\ve\dt\gamma p_{\min}^n\bnabla\!\cdot\!\bm u^{(1),n+1}-\ve^2\dt\gamma p_{\min}^n\bnabla\!\cdot\!\bm u^{(2),n+1}+{\cal O}(\ve^3).
\end{aligned}$$
Up to $\mathcal{O}(1)$ and $\mathcal{O}(\ve)$, we get $\bnabla\!\cdot\!\bm u^{(0),n+1}=0$ and $\bnabla\!\cdot\!\bm u^{(1),n+1}=0$ respectively.

Consistency with~\eqref{3.3} and~\eqref{3.4} comes directly from the structure of the scheme. In fact, thanks to what has been shown so far, up to $\mathcal{O}(1)$, the density and velocity updates~\eqref{4.11} and \eqref{4.12} give



%of the scheme \eref{4.8} with \eref{3.4} and \eref{3.5} as $\ve\to0$. We have already shown that
%$p^{(0),n+1}=p^{(0),n}\equiv\,\xbar p_0$ and $\bnabla\!\cdot\!\bm u^{(0),n+1}=0$. Finally, we rewrite \eref{4.15} as
%\begin{equation}
%\bm u^{(0),n+1}=\bm u^{(0),n}-\dt\bigg[(\bm u^{(0),n}\!\cdot\!\bnabla)\bm u^{(0),n}+\frac{\nabla p^{(2),n}}{\rho^{(0),n}}\bigg]
%-\frac{\dt}{\rho_{\max}^n}\left(\nabla p^{(2),n+1}-\nabla p^{(2),n}\right),
%\label{4.16}
%\end{equation}
%which is a consistent discretization of \eref{3.4}. 
\begin{align}
	&\begin{aligned}
		\rho^{(0),n+1}&=\rho^{(0),n}-\dt\bnabla\!\cdot\!\big(\rho^{(0),n}\bm u^{(0),n}\big),
	\end{aligned}\\
	&\begin{aligned}
		\bm u^{(0),n+1}&=\bm u^{(0),n}-\dt\bigg[(\bm u^{(0),n}\!\cdot\!\bnabla)\bm u^{(0),n}+\frac{\nabla p^{(2),n}}{\rho^{(0),n}}\bigg]-\frac{\dt}{\rho_{\max}^n}\left(\nabla p^{(2),n+1}-\nabla p^{(2),n}\right)\label{4.16},
	\end{aligned}
\end{align}
which are consistent with~\eqref{3.3} and~\eqref{3.4} respectively.

We remark that since the $(0)$- and $(1)$-modes of the pressure are constant, the evolution of the
pressure in the low-Mach-number limit essentially consists of the evolution of the $(2)$-mode. Thus,
$\nabla p^{(2),n+1}-\nabla p^{(2),n}\approx{\cal O}(\dt)$ and the last term in \eref{4.16}, in fact, represents a temporal diffusion term,
which is proportional to ${\cal O}((\dt)^2)$.
We also remark that, due to its definition in~\eqref{4.10}, the time step $\dt$ is asymptotically independent of $\ve$.
\hfill$\blacksquare$
\end{proof}

\subsubsection{Second-Order AP SI Time Discretization}\label{sec313}
We now introduce a second-order AP SI time discretization, which is based on the Deferred Correction (DeC) approach, which was \lore{originally introduced} 
%first proposed
in \cite{fox1949some}. Our second-order AP SI-DeC time discretization is directly related to the IMEX-DeC methods presented in
\cite{offner2025analysis} and based on the DeC formulation introduced in \cite{Decremi}, \lore{see also~\cite{micalizzi2025efficient,micalizzi2024new}}.

%According to the second-order AP SI-DeC time discretization, 
\lore{In particular, in the context of the adopted second-order AP SI-DeC time discretization,} the solution of \eref{4.47} is evolved from $t=t^n$ to $t=t^{n+1}$ \lore{through} the following two stages:
\begin{equation}
\begin{aligned}
&\bm V^*=\bm V^n-\dt\bm{{\cal R}}^n-\dt\bm{{\cal L}}^{n,*},\\
&\bm V^{n+1}=\bm V^n-\frac{\dt}{2}\big[\bm{{\cal R}}^n+\bm{{\cal R}}^*\big]-\frac{\dt}{2}\big[\bm{{\cal L}}^{n,n}-\bm{{\cal L}}^{*,*}\big]
-\dt\bm{{\cal L}}^{*,n+1},
\end{aligned}
\label{3.19}
\end{equation}
where the upper index $*$ \lore{is associated with} %denontes the data computed using 
the intermediate solution $\bm V^*$, and \lore{the} definitions of the operators
$\bm{{\cal R}}^*$ and $\bm{{\cal L}}^{n,n}$, $\bm{{\cal L}}^{*,*}$, and $\bm{{\cal L}}^{*,n+1}$ are \lore{analogous}
%similar 
to those given in \eref{3.12f} and \eref{3.13f}, respectively.

%The 
%System 
\lore{Scheme} \eref{3.19} can be equivalently written as
\begin{equation}
\begin{aligned}
\rho^*=&~\rho^n-\dt\bnabla\!\cdot\!(\rho^n\bm u^n),\\
\bm u^*=&~\bm u^n-\dt\Big[(\bm u^n\!\cdot\!\bnabla)\bm u^n+\frac{\rho_{\max}^n-\rho^n}{\ve^2\rho^n\rho_{\max}^n}\,\nabla p^n\Big]-
\frac{\dt}{\ve^2\rho_{\max}^n}\nabla p^*,\\
p^*=&~p^n-\dt\big[\bm u^n\!\cdot\!\nabla p^n+\gamma(p^n-p_{\min}^n)\bnabla\!\cdot\!\bm u^n\big]-\dt\gamma p_{\min}^n\bnabla\!\cdot\!\bm u^*,
\end{aligned}
\label{3.20}
\end{equation}
and
\begin{equation}
\begin{aligned}
\rho^{n+1}&=\rho^n-\frac{\dt}{2}\big[\bnabla\!\cdot\!(\rho^n\bm u^n)+\bnabla\!\cdot\!(\rho^*\bm u^*)\big],\\
\bm u^{n+1}&=\bm u^n-\frac{\dt}{2}\Big[(\bm u^n\!\cdot\!\bnabla)\bm u^n+(\bm u^*\!\cdot\!\bnabla)\bm u^*
+\frac{\rho_{\max}^n-\rho^n}{\ve^2\rho^n\rho_{\max}^n}\,\nabla p^n+\frac{\rho_{\max}^*-\rho^*}{\ve^2\rho^*\rho_{\max}^*}\,\nabla p^*\Big]\\
&\hspace*{1.05cm}-\frac{\dt}{2\ve^2}\left(\frac{\nabla p^n}{\rho^n_{\max}}-\frac{\nabla p^*}{\rho^*_{\max}}\right)-
\frac{\dt}{\ve^2\rho_{\max}^*}\nabla p^{n+1},\\[0.5ex]
p^{n+1}&=p^n-\frac{\dt}{2}\big[\bm u^n\!\cdot\!\nabla p^n+\bm u^*\!\cdot\!\nabla p^*+\gamma(p^n-p_{\min}^n)\bnabla\!\cdot\!\bm u^n
+\gamma(p^*-p_{\min}^*)\bnabla\!\cdot\!\bm u^*\big]\\[0.5ex]
&\hspace*{1.0cm}-\frac{\dt}{2}\,\gamma\big(p_{\min}^n\bnabla\!\cdot\!\bm u^n-p_{\min}^*\bnabla\!\cdot\!\bm u^*\big)
-\dt\gamma p_{\min}^*\bnabla\!\cdot\!\bm u^{n+1}.
\end{aligned}
\label{3.21}
\end{equation}

This second-order time discretization is indeed AP as shown in the next theorem.
\begin{theorem}
The second-order SI-DeC discretization \eref{3.20}--\eref{3.21} is AP according to Definition \ref{def41}, provided that $\dt$ is computed
according to \eref{4.10}.
\end{theorem}
\begin{proof}
We begin by observing that the first stage of the second-order SI-DeC discretization coincides with the first-order AP SI time
discretization studied above. Therefore, according to Theorem \ref{thm31} the intermediate solution $\rho^*$, $\bm u^*$, $p^*$, obtained by
\eref{3.20}, admits an expansion of the type \eref{4.1}, that is,
\begin{equation}
\rho^*=\rho^{(0),*}+\ve^2\rho^{(2),*}+\dots,\quad\bm u^*=\bm u^{(0),*}+\ve^2\bm u^{(2),*}+\dots,\quad p^*=p^{(0),*}+\ve^2p^{(2),*}+\dots,
\label{3.22}
\end{equation}
and \eref{3.20} is a consistent approximation of \eref{3.3}--\eref{3.5}, in particular,
\begin{equation}
\bnabla\!\cdot\!\bm u^{(0),*}=0,\quad p^{(0),*}\equiv\,\xbar p_0>0.
\label{3.23}
\end{equation}

We then substitute the expansions \eref{4.1} and \eref{3.22} into \eref{3.21} and use the conditions \eref{4.2} and \eref{3.23} to obtain
\begin{align}
\rho^{n+1}&=\rho^{(0),n}-\frac{\dt}{2}\Big[\bnabla\!\cdot\!\big(\rho^{(0),n}\bm u^{(0),n}\big)+
\bnabla\!\cdot\!\big(\rho^{(0),*}\bm u^{(0),*}\big)\Big]+{\cal O}(\ve^2),\label{3.24}\\
\bm u^{n+1}&=\bm u^{(0),n}-\frac{\dt}{2}\left[(\bm u^{(0),n}\!\cdot\!\bnabla)\bm u^{(0),n}+(\bm u^{(0),*}\!\cdot\!\bnabla)\bm u^{(0),*}+
\frac{\nabla p^{(2),n}}{\rho^{(0),n}}+\frac{\nabla p^{(2),*}}{\rho^{(0),*}}-\frac{2\nabla p^{(2),*}}{\rho_{\max}^*}\right]\nonumber\\
&-\frac{\dt}{\ve^2\rho_{\max}^*}\nabla p^{n+1}+{\cal O}(\ve^2),\label{3.25}\\[0.5ex]
p^{n+1}&=p^n-\dt\gamma p_{\min}^*\bnabla\!\cdot\!\bm u^{n+1}+{\cal O}(\ve^2).\label{3.26}
\end{align}
From \eref{3.24} we conclude that $\rho^{n+1}$ admits an expansion $\rho^{n+1}=\rho^{(0),n+1}+\ve^2\rho^{(2),n+1}+\dots$ with
$$
\rho^{(0),n+1}=\rho^{(0),n}-\frac{\dt}{2}\Big[\bnabla\!\cdot\!\big(\rho^{(0),n}\bm u^{(0),n}\big)+
\bnabla\!\cdot\!\big(\rho^{(0),*}\bm u^{(0),*}\big)\Big],
$$
which shows that the density equation in \eref{3.21} is consistent with \eref{3.3} as $\ve\to0$ since according to \eref{4.10},
$\dt$ is independent of $\ve$. Next, we take the divergence of the velocity equation \eref{3.25} and substitute $\bnabla\!\cdot\!\muu^{n+1}$
into the pressure equation \eref{3.26}. This yields
$$
\begin{aligned}
p^{n+1}-\frac{(\dt)^2\gamma p_{\min}^*}{\ve^2\rho_{\max}^*}\,\Delta p^{n+1}=
\,\xbar p_0&+\frac{(\dt)^2\gamma p_{\min}^*}{2}\,\bnabla\!\cdot\!\bigg[(\bm u^{(0),n}\!\cdot\!\bnabla)\bm u^{(0),n}+
(\bm u^{(0),*}\!\cdot\!\bnabla)\bm u^{(0),*}\\
&\hspace{3.3cm}+\frac{\nabla p^{(2),n}}{\rho^{(0),n}}+\frac{\nabla p^{(2),*}}{\rho^{(0),*}}-\frac{2\nabla p^{(2),*}}{\rho_{\max}^*}\bigg]+
{\cal O}(\ve^2),
\end{aligned}
$$
so that
$$
-\Delta p^{n+1}+\frac{\ve^2\rho_{\max}^*}{(\dt)^2\gamma p_{\min}^*}\,p^{n+1}={\cal O}(\ve^2),
$$
which implies $\Delta p^{n+1}={\cal O}(\ve^2)$. This together with the boundary conditions \eref{3.2} results in the same expansion for
$p^{n+1}$, which we have established in \eref{4.14} for the first-order SI method.

Next, we substitute \eref{4.14} into the the velocity equation \eref{3.25} to obtain
$$
\begin{aligned}
\bm u^{n+1}&=\bm u^{(0),n}-\frac{\dt}{2}\bigg[(\bm u^{(0),n}\!\cdot\!\bnabla)\bm u^{(0),n}+(\bm u^{(0),*}\!\cdot\!\bnabla)\bm u^{(0),*}+
\frac{\nabla p^{(2),n}}{\rho^{(0),n}}+\frac{\nabla p^{(2),*}}{\rho^{(0),*}}-\frac{2\nabla p^{(2),*}}{\rho_{\max}^*}\bigg]\\
&-\frac{\dt}{\rho_{\max}^*}\nabla p^{(2),n+1}+{\cal O}(\ve^2),
\end{aligned}
$$
which implies that $\bm u^{n+1}=\bm u^{(0),n+1}+\ve^2\bm u^{(2),n+1}+\dots$ with
\begin{equation}
\begin{aligned}
\bm u^{(0),n+1}&=\bm u^{(0),n}-\frac{\dt}{2}\bigg[(\bm u^{(0),n}\!\cdot\!\bnabla)\bm u^{(0),n}+(\bm u^{(0),*}\!\cdot\!\bnabla)\bm u^{(0),*}+
\frac{\nabla p^{(2),n}}{\rho^{(0),n}}+\frac{\nabla p^{(2),*}}{\rho^{(0),*}}-\frac{2\nabla p^{(2),*}}{\rho_{\max}^*}\bigg]\\
&-\frac{\dt}{\rho_{\max}^*}\nabla p^{(2),n+1}.
\end{aligned}
\label{4.15f}
\end{equation}

We now need to show the consistency of the scheme \eref{3.20}--\eref{3.21} with \eref{3.4} and \eref{3.5} as $\ve\to0$. We have already
shown that $p^{(0),n+1}=p^{(0),n}\equiv\,\xbar p_0$, which together with \eref{3.26} and the expansion of $\bm u^{n+1}$ implies
$\bnabla\!\cdot\!\bm u^{(0),n+1}\equiv0$. Finally, we rewrite \eref{4.15f} as
\begin{equation}
\begin{aligned}
\bm u^{(0),n+1}&=\bm u^{(0),n}-\frac{\dt}{2}\bigg[(\bm u^{(0),n}\!\cdot\!\bnabla)\bm u^{(0),n}+(\bm u^{(0),*}\!\cdot\!\bnabla)\bm u^{(0),*}+
\frac{\nabla p^{(2),n}}{\rho^{(0),n}}+\frac{\nabla p^{(2),*}}{\rho^{(0),*}}\bigg]\\
&-\frac{\dt}{\rho_{\max}^*}\left(\nabla p^{(2),n+1}-\nabla p^{(2),*}\right),
\end{aligned}
\label{4.16f}
\end{equation}
which is a consistent discretization of \eref{3.4}. We remark that as in \eref{4.16}, the last term in \eref{4.16f} represents the temporal
diffusion, which is consistent with the order of accuracy of the scheme and thus proportional to ${\cal O}((\dt)^3)$.\hfill$\blacksquare$
\end{proof}
\begin{remark}
The described AP SI-DeC time discretizations can be straightforwardly extended to arbitrarily high order within the DeC framework. For the
sake of brevity, we restrict our consideration to the second-order time accuracy, which matches the spatial accuracy of the developed AP
method.
\end{remark}

\subsection{\lore{Fully-Discrete} Second-Order AP Scheme for the Primitive System}\label{sec32}
In this section, we construct a \lore{fully-discrete} scheme based on the second-order AP SI time discretization presented in \S\ref{sec313}. To
this end, we first introduce uniform Cartesian cells $I_{j,k}:=[x_\jmh,x_\jph]\times[y_\kmh,y_\kph]$ with $x_\jph-x_\jmh\equiv\dx$ and
$y_\kph-y_\kmh\equiv\dy$, centered at $(x_j,y_k)$ with $x_j=\big(x_\jmh+x_\jph\big)/2$ and $y_k=\big(y_\kmh+y_\kph\big)/2$, and assume that
the cell averages $\,\xbar{\bm V}_{j,k}^{\,n}:\approx\iint_{I_{j,k}}\bm V(x,y,t^n)\,{\rm d}x{\rm d}y\,$ are available at time $t^n$.

A \lore{fully-discrete} FV verion of the second-order AP scheme \eref{3.19} reads as
\begin{align}
&\xbar{\bm V}^{\,*}_{j.k}=\,\xbar{\bm V}^{\,n}_{j.k}-\dt\bm{{\cal R}}^n_{j.k}-\dt\bm{{\cal L}}^{n,*}_{j.k},\label{3.31a}\\
&\xbar{\bm V}^{\,n+1}_{j.k}=\,\xbar{\bm V}^{\,n}_{j.k}-\frac{\dt}{2}\big[\bm{{\cal R}}^n_{j.k}+\bm{{\cal R}}^*_{j.k}\big]-
\frac{\dt}{2}\big[\bm{{\cal L}}^{n,n}_{j.k}-\bm{{\cal L}}^{*,*}_{j.k}\big]-\dt\bm{{\cal L}}^{*,n+1}_{j.k},\label{3.31b}
\end{align}
where $\bm{{\cal R}}^n_{j.k}$ and $\bm{{\cal R}}^*_{j.k}$ are obtained using the PCCU discretization from \cite{ACKM,CKN22}, which is a
low-dissipation generalization of the PCCU discretization from \cite{diaz2019path}, while $\bm{{\cal L}}^{n,*}_{j.k}$,
$\bm{{\cal L}}^{n,n}_{j.k}$, $\bm{{\cal L}}^{*,*}_{j.k}$, and $\bm{{\cal L}}^{*,n+1}_{j.k}$ are obtained using centaral differences. In what
follows, we provide details on $\bm{{\cal R}}^n_{j.k}$ and $\bm{{\cal L}}^{n,*}_{j.k}$ only, whereas the remaining discretizations are
obtained in a similar manner.

We begin with
\begin{equation}
\begin{aligned}
\bm{{\cal R}}_{j,k}^n=&\,\frac{1}{\dx}\Bigg[\widetilde{\bm{{\cal F}}}_{\jph,k}^n-\widetilde{\bm{{\cal F}}}_{\jmh,k}^n-
\widetilde{\bm B}_{j,k}^n-\frac{\tilde a^{+,n}_{\jmh,k}\widetilde{\bm B}_{\bm\Psi,\jmh,k}^n}
{\tilde a^{+,n}_{\jmh,k}-\tilde a^{-,n}_{\jmh,k}}+
\frac{\tilde a^{-,n}_{\jph,k}\widetilde{\bm B}_{\bm\Psi,\jph,k}^n}{\tilde a^{+,n}_{\jph,k}-\tilde a^{-,n}_{\jph,k}}\Bigg]\\
+&\,\frac{1}{\dy}\Bigg[\widetilde{\bm{{\cal G}}}_{j,\kph}^n-\widetilde{\bm{{\cal G}}}_{j,\kmh}^n-\widetilde{\bm C}_{j,k}^n
-\frac{\tilde b^{+,n}_{j,\kmh}\widetilde{\bm C}_{\bm\Psi,j,\kmh}^n}{\tilde b^{+,n}_{j,\kmh}-\tilde b^{-,n}_{j,\kmh}}+
\frac{\tilde b^{-,n}_{j,\kph}\widetilde{\bm C}_{\bm\Psi,j,\kph}^n}{\tilde b^{+,n}_{j,\kph}-\tilde b^{-,n}_{j,\kph}}\Bigg],
\end{aligned}
\label{3.32}
\end{equation}
where $\widetilde{\bm{{\cal F}}}_{\jph,k}^n$ and $\widetilde{\bm{{\cal G}}}_{j,\kph}^n$ are the CU numerical fluxes
$$
\begin{aligned}
&\widetilde{\bm{{\cal F}}}_{\jph,k}^n=\frac{\tilde a^{+,n}_{\jph,k}\widetilde{\bm F}\big(\bm V^{-,n}_{\jph,k}\big)-
\tilde a^{-,n}_{\jph,k}\widetilde{\bm F}\big(\bm V^{+,n}_{\jph,k}\big)}{\tilde a^{+,n}_{\jph,k}-\tilde a^{+,n}_{\jph,k}}+
\frac{\tilde a^{+,n}_{\jph,k}\tilde a^{-,n}_{\jph,k}}{\tilde a^{+,n}_{\jph,k}-\tilde a^{-,n}_{\jph,k}\,}
\Big(\bm V^{+,n}_{\jph,k}-\bm V^{-,n}_{\jph,k}-\delta\bm V_{\jph,k}^n\Big),\\
&\widetilde{\bm{{\cal G}}}_{j,\kph}^n=\frac{\tilde b^{+,n}_{j,\kph}\widetilde{\bm G}\big(\bm V^{-,n}_{j,\kph}\big)-
\tilde b^{-,n}_{j,\kph}\widetilde{\bm G}\big(\bm V^{+,n}_{j,\kph}\big)}{\tilde b^{+,n}_{j,\kph}-\tilde b^{-,n}_{j,\kph}}+
\frac{\tilde b^{+,n}_{j,\kph}\tilde b^{-,n}_{j,\kph}}{\tilde b^{+,n}_{j,\kph}-\tilde b^{-,n}_{j,\kph}}\,
\Big(\bm V^{+,n}_{j,\kph}-\bm V^{-,n}_{j,\kph}-\delta\bm V_{j,\kph}^n\Big),
\end{aligned}
$$
where $\bm V^{\pm,n}_{\jph,k}$ and $\bm V^{\pm,n}_{j,\kph}$ are point values of $\bm V$ at the midpoints of the cell $I_{j,k}$,
$\delta\bm V_{\jph,k}^n$ and $\delta\bm V_{j,\kph}^n$ are ``built-in'' anti-diffusion terms, and $\tilde a^{\pm,n}_{\jph,k}$ and
$\tilde a^{\pm,n}_{j,\kph}$ denote the one-sided local propagation speeds of in the $x$- and $y$-directions, respectively.

The point values
\begin{equation}
\begin{aligned}
&\bm V^{-,n}_{\jph,k}=\,\xbar{\bm V}_{j,k}^{\,n}+\frac{\dx}{2}(\bm V_x)_{j,k}^n,&&
\bm V^{+,n}_{\jph,k}=\,\xbar{\bm V}_{j+1,k}^{\,n}-\frac{\dx}{2}(\bm V_x)_{j+1,k}^n,\\
&\bm V^{-,n}_{j,\kph}=\,\xbar{\bm V}_{j,k}^{\,n}+\frac{\dy}{2}(\bm V_y)_{j,k}^n,&&
\bm V^{+,n}_{j,\kph}=\,\xbar{\bm V}_{j,k+1}^{\,n}-\frac{\dy}{2}(\bm V_y)_{j,k+1}^n,
\end{aligned}
\label{3.33f}
\end{equation}
are computed using the piecewise linear reconstruction
$$
\xbar{\bm V}_{j,k}^{\,n}+(\bm V_x)_{j,k}^n(x-x_j)+(\bm V_y)_{j,k}^n(y-y_k),~~(x,y)\in I_{j,k},
$$
in which the slopes $(\bm V_x)_{j,k}^n$ and $(\bm V_y)_{j,k}^n$ are approximated using the generalized minmod limiter (see, e.g.,
\cite{LieNoe,nessyahu1990non,sweby1984high}):
\begin{equation}
\begin{aligned}
(\bV_x)_{j,k}^n&={\rm minmod}\left(\theta\,\frac{\xbar{\bV}_{j,k}^{\,n}-\,\xbar{\bV}_{j-1,k}^{\,n}}{\dx},\,
\frac{\xbar{\bV}_{j+1,k}^{\,n}-\,\xbar{\bV}_{j-1,k}^{\,n}}{2\dx},\,
\theta\,\frac{\xbar{\bV}_{j+1,k}^{\,n}-\,\xbar{\bV}_{j,k}^{\,n}}{\dx}\right),\\
(\bV_y)_{j,k}^n&={\rm minmod}\left(\theta\,\frac{\xbar{\bV}_{j,k}^{\,n}-\,\xbar{\bV}_{j,k-1}^{\,n}}{\dy},\,
\frac{\xbar{\bV}_{j,k+1}^{\,n}-\,\xbar{\bV}_{j,k-1}^{\,n}}{2\dy},\,
\theta\,\frac{\xbar{\bV}_{j,k+1}^{\,n}-\,\xbar{\bV}_{j,k}^{\,n}}{\dy}\right),
\end{aligned}
\label{3.33}
\end{equation}
where the minmod function, defined by
\begin{equation*}
{\rm minmod}(z_1,z_2,\ldots)=\left\{\begin{aligned}
&\min(z_1,z_2,\ldots)&&\mbox{if}~z_i>0,~\forall i,\\
&\max(z_1,z_2,\ldots)&&\mbox{if}~z_i<0,~\forall i,\\
&\,0&&\mbox{otherwise},
\end{aligned}\right.
\end{equation*}
is applied in a componentwise manner. The parameter $\theta\in[1,2]$ in \eref{3.33} is to be chosen to adjust the amount of numerical
dissipation present in the resulting scheme, with larger values of $\theta$ leading to sharper but, in general, more oscillatory solutions.

The one-sided local speeds of propagation are then estimated using the smallest and largest eigenvalues of matrices
$\frac{\partial\widetilde{\mF}}{\partial\bm V}(\bm V)-\widetilde B(\bm V)$ and
$\frac{\partial\widetilde{\mG}}{\partial\bm V}(\bm V)-\widetilde C(\bm V)$ as follows:
\begin{equation}
\begin{aligned}
&\tilde a^{-,n}_{\jph,k}=\min\left\{u^{-,n}_{\jph,k}-\tilde c^{\,-,n}_{\jph,k},\,u^{+,n}_{\jph,k}-\tilde c^{\,+,n}_{\jph,k},\,-\delta
\right\},\\
&\tilde a^{+,n}_{\jph,k}=\max\left\{u^{-,n}_{\jph,k}+\tilde c^{\,-,n}_{\jph,k},\,u^{+,n}_{\jph,k}+\tilde c^{\,+,n}_{\jph,k},\,
\delta\right\},\\
&\tilde b^{-,n}_{j,\kph}=\min\left\{v^{-,n}_{j,\kph}-\tilde c^{\,-,n}_{j,\kph},\,v^{+,n}_{j,\kph}-\tilde c^{\,+,n}_{j,\kph},\,
-\delta\right\},\\
&\tilde b^{+,n}_{j,\kph}=\max\left\{v^{-,n}_{j,\kph}+\tilde c^{\,-,n}_{j,\kph},\,v^{+,n}_{j,\kph}+\tilde c^{\,+,n}_{j,\kph},\,
\delta\right\},
\end{aligned}
\label{4.49}
\end{equation}
where the sound speeds
$$
\tilde c^{\,\pm,n}_{\jph,k}=\frac{1}{\ve}\sqrt{\gamma\frac{\big(\rho_{\max}^n-\rho^{\,\pm,n}_{\jph,k}\big)
\big(p^{\,\pm,n}_{\jph,k}-p_{\min}^n)}{\rho^{\,\pm,n}_{\jph,k}\,\rho_{\max}^n}},\quad
\tilde c^{\,\pm,n}_{j,\kph}=\frac{1}{\ve}\sqrt{\gamma\frac{\big(\rho_{\max}^n-\rho^{\,\pm,n}_{j,\kph}\big)
\big(p^{\,\pm,n}_{j,\kph}-p_{\min}^n)}{\rho^{\,\pm,n}_{j,\kph}\,\rho_{\max}^n}}
$$
are computed using the following discrete versions of $\rho_{\max}^n$ and $p_{\min}^n$:
\begin{equation}
\rho_{\max}^n=\max\limits_{j,k}\,\xbar{\rho}_{j,k}^{\,n},\quad p_{\min}^n=\min\limits_{j,k}\,\xbar p_{j,k}^{\,n},
\label{3.35}
\end{equation}
and $\delta$ is a small positive parameter introduced to prevent divisions by $0$ (we have taken $\delta=10^{-15}$ in the numerical
experiments reported in \S\ref{sec6}).

The ``built-in'' anti-diffusion terms are 
$$
\begin{aligned}
&\delta\bm V_{\jph,k}^n={\rm minmod}\Big(\bm V_{\jph,k}^{{\rm int},n}-\bm V^{-,n}_{\jph,k},\,
\bm V^{+,n}_{\jph,k}-\bm V_{\jph,k}^{{\rm int},n}\Big),\\
&\delta\bm V_{j,\kph}^n={\rm minmod}\Big(\bm V_{j,\kph}^{{\rm int},n}-\bm V^{-,n}_{j,\kph},\,
\bm V^{+,n}_{j,\kph}-\bm V_{j,\kph}^{{\rm int},n}\Big),
\end{aligned}
$$
where 
$$
\begin{aligned}
&\bm V_{\jph,k}^{{\rm int},n}=\frac{\tilde a^{+,n}_{\jph,k}\bm V^{+,n}_{\jph,k}-\tilde a^{-,n}_{\jph,k}\bm V^{-,n}_{\jph,k}
-\widetilde{\bm F}\big(\bm V^{+,n}_{\jph,k}\big)+\widetilde{\bm F}\big(\bm V^{-,n}_{\jph,k}\big)}
{\tilde a^{+,n}_{\jph,k}-\tilde a^{-,n}_{\jph,k}},\\[0.3ex]
&\bm V_{j,\kph}^{{\rm int},n}=\frac{\tilde b^{+,n}_{j,\kph}\bm V^{+,n}_{j,\kph}-\tilde b^{-,n}_{j,\kph}\bm V^{-,n}_{j,\kph}
-\widetilde{\bm G}\big(\bm V^{+,n}_{j,\kph}\big)+\widetilde{\bm G}\big(\bm V^{-,n}_{j,\kph}\big)}
{\tilde b^{+,n}_{j,\kph}-\tilde b^{-,n}_{j,\kph}}.
\end{aligned}
$$

Finally,
\begin{equation}
\bm{{\cal L}}^{n,*}_{j.k}=\Big(0,\,\frac{\,\xbar p_{j+1,k}^{\,*}-\,\xbar p_{j-1,k}^{\,*}}{2\dx\,\ve^2\rho_{\max}^n},\,
\frac{\,\xbar p_{j,k+1}^{\,*}-\,\xbar p_{j,k-1}^{\,*}}{2\dy\,\ve^2\rho_{\max}^n},\,
\gamma p_{\min}^n\bnabla\!\cdot\xbar{\bm u}_{j,k}^{\,*}\Big)^\top,
\label{4.48b}
\end{equation}
where $\bnabla\!\cdot\xbar{\bm u}_{j,k}$ denote the discrete divergence operator computed using the second-order central differences:
\begin{equation}
\bnabla\!\cdot\xbar{\bm u}_{j,k}:=\frac{\,\xbar u_{j+1,k}-\,\xbar u_{j-1,k}}{2\dx}+\frac{\,\xbar v_{j,k+1}-\,\xbar v_{j,k-1}}{2\dy}.
\label{3.38f}
\end{equation}

\subsection{Semi-Discrete CU Scheme for the Conservative System}\label{sec42}
We now consider the conservative formulation \eref{2.1a}--\eref{2.1c}, which can be into the following vector form:
\begin{equation}
\begin{aligned}
&\bm U_t+\mF(\bm U)_x+\mG(\bm U)_y=\bm0,\quad\bm U:=(\rho,\rho u,\rho v,E)^\top,\\
&\mF(\bm U):=\Big(\rho u,\rho u^2+\frac{p}{\ve^2},\rho uv,u(E+p)\Big)^\top,\quad
\mG(\bm U):=\Big(\rho v,\rho uv,\rho v^2+\frac{p}{\ve^2},v(E+p)\Big)^\top.
\end{aligned}
\label{3.39g}
\end{equation}

In the semi-discrete CU scheme, the cell averages $\,\xbar{\bm U}_{j,k}(t):\approx\iint_{I_{j,k}}\bm U(x,y,t)\,{\rm d}x{\rm d}y$ are evolved
in time by numerically solving the following system of ODEs:
\begin{equation}
\frac{{\rm d}}{{\rm d}t}\,\xbar{\bm U}_{j,k}=-\frac{\bm{{\cal F}}_{\jph,k}-\bm{{\cal F}}_{\jmh,k}}{\dx}-
\frac{\bm{{\cal G}}_{j,\kph}-\bm{{\cal G}}_{j,\kmh}}{\dy},
\label{2.15a}
\end{equation}
where $\bm{{\cal F}}_{\jph,k}$ and $\bm{{\cal G}}_{j,\kph}$ are the CU numerical fluxes from \cite{KLin} defined as
\begin{equation}
\begin{aligned}
\bm{{\cal F}}_{\jph,k}&=\frac{a^+_{\jph,k}\bm F\big(\bm U^-_{\jph,k}\big)-a^-_{\jph,k}\bm F\big(\bm U^+_{\jph,k}\big)}
{a^+_{\jph,k}-a^-_{\jph,k}}\\
&+\frac{a^+_{\jph,k}a^-_{\jph,k}}{a^+_{\jph,k}-a^-_{\jph,k}}\,\Big(\bm U^+_{\jph,k}-\bm U^-_{\jph,k}-\delta\bm U_{\jph,k}\Big),\\[0.8ex]
\bm{{\cal G}}_{j,\kph}&=\frac{b^+_{j,\kph}\bm G\big(\bm U^-_{j,\kph}\big)-b^-_{j,\kph}\bm G\big(\bm U^+_{j,\kph}\big)}
{b^+_{j,\kph}-b^-_{j,\kph}}\\
&+\frac{b^+_{j,\kph}b^-_{j,\kph}}{b^+_{j,\kph}-b^-_{j,\kph}}\,\Big(\bm U^+_{j,\kph}-\bm U^-_{j,\kph}-\delta\bm U_{j,\kph}\Big).
\end{aligned}
\label{3.41}
\end{equation}
Here, the interface values $\bm U_{\jph,k}^\pm:=\bm U\big(\bm V_{\jph,k}^\pm\big)$ and
$\bm U_{j,\kph}^\pm:=\bm U\big(\bm V_{j,\kph}^\pm\big)$ are computed from the reconstructed primitive variables $\bm V_{\jph,k}^\pm$ and 
$\bm V_{j,\kph}^\pm$ (see \S\ref{sec32}) at the corresponding time level via a straightforward transformation $\bU(\bV)$ from $\bm V$ to
$\bm U$. In addition, $a_{\jph,k}^\pm$ and $b_{j,\kph}^\pm$ are the one-sided local speeds of propagation for the conservative system
\eref{3.39g} in the $x$- and $y$-directions, respectively. They are estimated using the largest and 
smallest eigenvalues of the corresponding flux Jacobians as follows:
\begin{equation}
\begin{aligned}
\begin{aligned}
&a^-_{\jph,k}=\min\left\{u^-_{\jph,k}-c^-_{\jph,k},\,u^+_{\jph,k}-c^+_{\jph,k},\,-\delta\right\},\\
&a^+_{\jph,k}=\max\left\{u^-_{\jph,k}+c^-_{\jph,k},\,u^+_{\jph,k}+c^+_{\jph,k},\,\delta\right\},
\end{aligned}\quad c^\pm_{\jph,k}=\frac{1}{\ve}\sqrt{\frac{\gamma p^\pm_{\jph,k}}{\rho^\pm_{\jph,k}}},\\[1.5ex]
\begin{aligned}
&b^-_{j,\kph}=\min\left\{v^-_{j,\kph}-c^-_{j,\kph},\,v^+_{j,\kph}-c^+_{j,\kph},\,-\delta\right\},\\
&b^+_{j,\kph}=\max\left\{v^-_{j,\kph}+c^-_{j,\kph},\,v^+_{j,\kph}+c^+_{j,\kph},\,\delta\right\},
\end{aligned}\quad c^\pm_{j,\kph}=\frac{1}{\ve}\sqrt{\frac{\gamma p^\pm_{j,\kph}}{\rho^\pm_{j,\kph}}},
\end{aligned}
\label{3.42}
\end{equation}
where $\delta=10^{-15}$ is used to avoid divisions by $0$.

The ``built-in'' anti-diffusion terms are 
\begin{equation}
\begin{aligned}
&\delta\bm U_{\jph,k}={\rm minmod}\Big(\bm U_{\jph,k}^{\rm int}-\bm U^-_{\jph,k},\bm U^+_{\jph,k}-\bm U_{\jph,k}^{\rm int}\Big),\\[0.5ex]
&\delta\bm U_{j,\kph}={\rm minmod}\Big(\bm U_{j,\kph}^{\rm int}-\bm U^-_{j,\kph},\bm U^+_{j,\kph}-\bm U_{j,\kph}^{\rm int}\Big)
\end{aligned}
\label{3.43}
\end{equation}
with 
\begin{equation}
\begin{aligned}
&\bm U_{\jph,k}^{\rm int}=\frac{a^+_{\jph,k}\bm U^+_{\jph,k}-a^-_{\jph,k}\bm U^-_{\jph,k}
-\bm F(\bm U^+_{\jph,k})+\bm F(\bm U^-_{\jph,k})}{a^+_{\jph,k}-a^-_{\jph,k}},\\[0.5ex]
&\bm U_{j,\kph}^{\rm int}=\frac{b^+_{j,\kph}\bm U^+_{j,\kph}-b^-_{j,\kph}\bm U^-_{j,\kph}
-\bm G\big(\bm U^+_{j,\kph}\big)+\bm G\big(\bm U^-_{j,\kph}\big)}{b^+_{j,\kph}-b^-_{j,\kph}}.
\end{aligned}
\label{3.44}
\end{equation}

Note that most of the indexed quantities in the semi-discrete setting above are time-dependent, but we have omitted this dependence for the
sake of brevity.

Finally, the system of ODEs \eref{2.15a} has to be integrated in time using an appropriate ODE solver. We take advantage of the DF-FV 
approach, which allows to use an implicit trapezoidal method as explained in \S\ref{sec422}.

\subsection{Implementation Details}\label{sec422}
In the DF-FV approach, the solutions of the primitive and conservative systems are evolved simultaneously according to the following
algorithm.

First, we update the solution of the SI scheme \eref{3.31a}--\eref{3.31b} in the following six steps.

\smallskip
\noindent
$\bullet\,$ \textbf{Step 1 (Compute $\xbar\rho^{\,*}_{j,k}$).} We use the $\rho$-equation in \eref{3.31a} to obtain
\begin{equation*}
\xbar\rho^{\,*}_{j,k}=\,\xbar\rho^{\,n}_{j,k}-\dt({\cal R}^\rho)^n_{j,k}.
\end{equation*}

\smallskip
\noindent
$\bullet\,$ \textbf{Step 2 (Solve the linear elliptic equation for $\,\xbar p^{\,*}_{j,k}$).} We apply the discrete divergence operator
\eref{3.38f} to the $\bm u$-equations in \eref{3.31a} and substitute them into the $p$-equation in \eref{3.31a} to obtain the following
linear system of algebraic equations for $\,\xbar p^{\,*}_{j,k}$, which is a discretization of the linear elliptic equation for $p^*$:
\begin{equation*}
\xbar p^{\,*}_{j,k}-\frac{(\dt)^2\gamma p_{\min}^n}{\ve^2\rho_{\max}^n}\,\Delta\xbar p_{j,k}^{\,*}=\,\xbar p^{\,n}_{j,k}-
\dt({\cal R}^p)^n_{j,k}-\dt\gamma p_{\min}^n\bnabla\!\cdot\xbar{\bm u}_{j,k}^{\,n}+
+(\dt)^2\gamma p_{\min}^n\bnabla\!\cdot\!(\bm{{\cal R}}^{\bm u})^n_{j,k},
\label{3.39f}
\end{equation*}
where the discrete Laplacian $\Delta\xbar p_{j,k}$ is defined as
$$
\Delta\xbar p_{j,k}:=\frac{\,\xbar p_{j-1,k}-2\,\xbar p_{j,k}+\,\xbar p_{j+1,k}}{(\dx)^2}+
\frac{\,\xbar p_{j,k-1}-2\,\xbar p_{j,k}+\,\xbar p_{j,k+1}}{(\dy)^2}.
$$

\smallskip
\noindent
$\bullet\,$ \textbf{Step 3 (Compute $\,\xbar{\bm u}^{\,*}_{j,k}$).} Once $\,\xbar p^{\,*}_{j,k}$ is available, we use the $\bm u$-equations
in \eref{3.31a} to obtain
\begin{equation*}
\xbar{\bm u}^{\,*}_{j,k}=\,\xbar{\bm u}^{\,n}_{j,k}-\dt(\bm{{\cal R}}^{\bm u})^n_{j,k}-\dt(\bm{{\cal L}}^{\bm u})^{n,*}_{j,k}.
\end{equation*}

\smallskip
\noindent
$\bullet\,$ \textbf{Step 4 (Compute $\xbar\rho^{\,n+1}_{j,k}$).} We solve the $\rho$-equation in \eref{3.31b} to obtain
\begin{equation*}
\xbar\rho^{\,n+1}_{j,k}=\,\xbar\rho^{\,n}_{j,k}-\frac{\dt}{2}\big[({\cal R}^\rho)^n_{j,k}+({\cal R}^\rho)^*_{j,k}\big].
\end{equation*}

\smallskip
\noindent
$\bullet\,$ \textbf{Step 5 (Solve the linear elliptic equation for $\,\xbar p^{\,n+1}_{j,k}$).} We apply the discrete divergence operator
\eref{3.38f} to the $\bm u$-equations in \eref{3.31b} and substitute them into the $p$-equation in \eref{3.31b} to obtain the following
linear system of algebraic equations for $\,\xbar p^{\,n+1}_{j,k}$, which is a discretization of the linear elliptic equation for $p^{n+1}$:
\begin{equation*}
\begin{aligned}
\xbar p^{\,n+1}_{j,k}&-\frac{(\dt)^2\gamma p_{\min}^*}{\ve^2\rho_{\max}^*}\,\Delta\xbar p_{j,k}^{\,n+1}=\,\xbar p^{\,n}_{j,k}-
\frac{\dt}{2}\left[({\cal R}^p)_{j,k}^n+({\cal R}^p)_{j,k}^*\right]-
\frac{\dt}{2}\left[({\cal L}^p)^{n,n}_{j,k}-({\cal L}^p)^{*,*}_{j,k}\right]\\
&-\dt\gamma p_{\min}^*\bnabla\!\cdot\xbar{\bm u}_{j,k}^{\,n}
+\frac{(\dt)^2\gamma p_{\min}^*}{2}\bnabla\!\cdot\!\left[(\bm{{\cal R}}^{\bm u})^n_{j,k}+(\bm{{\cal R}}^{\bm u})^*_{j,k}\right]
-\frac{(\dt)^2\gamma p_{\min}^*}{2}\bnabla\!\cdot\!\left[(\bm{{\cal L}}^u)^{n,n}_{j,k}-(\bm{{\cal L}}^u)^{*,*}_{j,k}\right].
\end{aligned}
\label{3.40f}
\end{equation*}

\smallskip
\noindent
$\bullet\,$ \textbf{Step 6 (Compute $\,\xbar{\bm u}^{\,n+1}_{j,k}$).} Once $\,\xbar p^{\,n+1}_{j,k}$ is available, we compute
\begin{equation*}
\begin{aligned}
\xbar{\bm u}^{\,n+1}_{j,k}=
\,\xbar{\bm u}^{\,n}_{j,k}-\frac{\dt}{2}\left[(\bm{{\cal R}}^{\bm u})^n_{j,k}+(\bm{{\cal R}}^{\bm u})^*_{j,k}\right]
-\frac{\dt}{2}\left[(\bm{{\cal L}}^u)^{n,n}_{j,k}-(\bm{{\cal L}}^u)^{*,*}_{j,k}\right]-\dt(\bm{{\cal L}}^{\bm u})^{*,n+1}_{j,k}.
\end{aligned}
\end{equation*}

Second, equipped with both $\,\xbar{\bm V}^{\,n}_{j,k}$ and $\,\xbar{\bm V}^{\,n+1}_{j,k}$, we proceed with solving the ODE system
\eref{2.15a} using the thapezoidal method.

\smallskip
\noindent
$\bullet\,$ \textbf{Step 7 (Compute $\,\xbar{\bm U}^{\,n+1}_{j,k}$).} We evaluate the numerical fluxes \eref{3.41}--\eref{3.44} at both the
time levels $t^n$ and $t^{n+1}$ and then obtain
\begin{equation}
\begin{aligned}
\xbar{\bm U}^{\,n+1}_{j,k}=\,\xbar{\bm U}^{\,n}_{j,k}&-\frac{\dt}{2}\left[\frac{\bm{{\cal F}}_{\jph,k}^n-\bm{{\cal F}}_{\jmh,k}^n}{\dx}+
\frac{\bm{{\cal F}}_{\jph,k}^{n+1}-\bm{{\cal F}}_{\jmh,k}^{n+1}}{\dx}\right]\\
&-\frac{\dt}{2}\left[\frac{\bm{{\cal G}}_{j,\kph}^n-\bm{{\cal G}}_{j,\kmh}^n}{\dy}+
\frac{\bm{{\cal G}}_{j,\kph}^{n+1}-\bm{{\cal G}}_{j,\kmh}^{n+1}}{\dy}\right].
\end{aligned}
\label{3.45}
\end{equation}
Recall that the interface values of $\bm U$ needed for the computation of the numerical fluxes in \eref{3.45} are obtained from the
reconstructed primitive variables $\bm V$ at the corresponding time levels. Note that the conservative update \eref{3.45} is, as a matter of
fact, explicit, since the computation of $\xbar{\bm V}^{\,n+1}_{j,k}$ is performed independently of the conserved system.

\subsection{Post-Processing}\label{sec43}
Upon the completion of Step 7, we obtain two sets of numerical solutions at time $t=t^{n+1}$: the $\bm V$-solution, which is AP but  
nonconservative, and the $\bm U$-solution, which is conservative but non-AP. We therefore design a post-processing \eref{4.5f} using their
convex combination with coefficients proportional to $\ve$, leveraging the AP SI method in the low-Mach-number regime and the sharp
conservative CU scheme in the moderate- and high-Mach-number regimes---thus ensuring accuracy, stability, and physical consistency across
all flow regimes. Specifically, we select the following function $r$ in \eref{4.5f}:
\begin{equation}
r\Big(\xbar{\bm V}\big(\,\xbar{\bm U}^{\,n+1}_{j,k}\big),\,\xbar{\bm V}^{\,n+1}_{j,k}\Big)=
(1-s(\ve))\,\xbar{\bm V}\big(\,\xbar{\bm U}^{\,n+1}_{j,k}\big)+s(\ve)\,\xbar{\bm V}^{\,n+1}_{j,k},\quad\forall j,k,
\label{3.46}
\end{equation}
where $s$ is a suitable switching function, which is supposed to be increasing, continuous, and satisfy $s(1)=0$ and $s(0)=1$. Moreover,
in the high-Mach-number regime, $s$ should be $\sim0$ so that the primitive variables $\bm V^{n+1}$ are almost completely overwritten by
$\bm V(\bm U^{n+1})$, while in the low-Mach-number regime, $s$ should be $\sim1$ so that the primitive variables $\bm V^{n+1}$ stay almost
unchanged. For intermediate values of $\ve$, a smooth transition between $1$ and $0$ is expected.

\section{Numerical Examples}\label{sec6}
In this section, we verify the accuracy and robustness of the proposed AP scheme on a variety of numerical examples across different values 
of $\ve$. In all of the numerical examples, we:

\smallskip
\noindent
$\bullet$ Take the minmod parameter $\theta=1.3$;

\smallskip
\noindent
$\bullet$ Adaptively select time steps based on the time-step restriction \eref{4.10} for the nonstiff part of the primitive system. 

\smallskip
\noindent
$\bullet$ Set $\gamma=1.4$ (except for Example 1, in which $\gamma=2$);

\smallskip
\noindent
$\bullet$ Modify \eref{4.3} to
\begin{equation}
\rho_{\max}=\max\limits_{(x,y)\in\Omega}\rho+\ve^4,\quad p_{\min}=\min\limits_{(x,y)\in\Omega}p-\ve^4,
\label{4.9}
\end{equation}
Notice that this modification has almost no impact in the low-Mach-number regime, but it aims at adding more upwinding and thus improving
the stability property of the resulting AP scheme when $\ve$ is large;

\smallskip
\noindent
$\bullet$ Choose the following switching function:
\begin{equation}
s(\ve)=\begin{cases}
1-\ve^{14},&0<\ve\le0.15,\\
\exp\bigg(1-\dfrac{1}{1-16\left(\ve-0.15\right)^2}\bigg)\left[0.85^{14}-0.6^{14}\right]+0.6^{14},&0.15<\ve<0.4,\\
(1-\ve)^{14},&0.4\le\ve\le1,
\end{cases}
\label{3.47}
\end{equation}
which is plotted in Figure \ref{fig31}.
\begin{figure}[ht!]
\centerline{\includegraphics[width=0.32\textwidth]{switching_function}}
\caption{\sf Switching function $s(\varepsilon)$ plotted with respect to $1/\ve$.\label{fig31}}
\end{figure}

\subsubsection*{Example 1---Accuracy Test for Low-Mach-Number Smooth Vortex}
In this example taken from \cite{zeifang2020novel}, we consider a smooth unsteady Mach dependent vortex over the computational domain
$[-10,10]\times[-10,10]$ subject to the periodic boundary conditions. The analytical solution is given, modulo the periodicity, by
\begin{equation*}
\begin{aligned}
\rho(\bm x_r)&=1-\frac{\ve^2}{16\pi^2}\,{\rm e}^{1-\|\bm x_r\|_2^2},\quad
u(\bm x_r)=1-\frac{\ve y_r}{2\pi}\,{\rm e}^{\frac{1-\|\bm x_r\|_2^2}{2}},\quad
v(\bm x_r)=1+\frac{\ve x_r}{2\pi}\,{\rm e}^{\frac{1-\|\bm x_r\|_2^2}{2}},\\
E(\bm x_r)&=1+\ve^2\Big[\rho^2(\bm x_r)+\frac{\rho(\bm x_r)}{2}\big(u^2(\bm x_r)+v^2(\bm x_r)\big)\Big],
\end{aligned}
\end{equation*}
where $\bm x_r(x,y,t)=(x_r,y_r)^\top:=(x-t,y-t)^\top$.

We take the CFL number $K_{\rm CFL}=0.475$ and compute the numerical solution until the final time $t=0.1$ on a series of uniform
$N\times N$ meshes with $N=64$, $128$, $256$, and $512$ for $\ve=1$, $0.1$, $0.01$, and $0.001$. We study the convergence and present the
obtained results in Figure \ref{fig42}, where one can see that the expected second-order convergence rate has been achieved in all variables
for all considered $\ve$. One can also observe that for fixed mesh refinement, the error decreases for decreasing $\ve$ as a result of the
convergence of the analytical solution to the incompressible limit $(\rho,u,v,p)^\top=(1,1,1,1)^\top$ and of the AP character of the
proposed AP DF-FV scheme.
\begin{figure}[ht!]
\centerline{\includegraphics[width=0.28\textwidth]{AP_paper_convergence_low_mach_vortex_64_512_CFL0.475_rho}\hspace*{0.8cm}
	    \includegraphics[width=0.28\textwidth]{AP_paper_convergence_low_mach_vortex_64_512_CFL0.475_p}}
\vskip10pt
\centerline{\includegraphics[width=0.28\textwidth]{AP_paper_convergence_low_mach_vortex_64_512_CFL0.475_u}\hspace*{0.8cm}
	    \includegraphics[width=0.28\textwidth]{AP_paper_convergence_low_mach_vortex_64_512_CFL0.475_v}}
\caption{\sf Example 1: Convergence analysis.\label{fig42}}
\end{figure}

\subsubsection*{Example 2---Gresho Vortex}
This example was introduced in \cite{gresho1990theory} and since then it has been widely used as a common benchmark to numerically validate
the AP property. We consider a steady vortex over the computational domain $[0,1]\times[0,1]$ subject to the periodic boundary conditions.
At any time $t$, the shape of the vortex is given by
\begin{equation*}
\begin{aligned}
\rho(r)&\equiv1,\quad u(r)=-\frac{y_r}{r}\psi(r),\quad v(r)=\frac{x_r}{r}\psi(r),\\
p(r)&=\begin{cases}
1+12.5\ve^2r^2,&r<0.2,\\
1+\ve^2\left(4\ln(5r)+4-20r+12.5r^2\right),&0.2\le r<0.4,\\
1+\ve^2(4\ln2-2),&r\ge0.4,
\end{cases}
\end{aligned}
\end{equation*}
where
\begin{equation*}
x_r:=x-0.5,\quad y_r:=y-0.5,\quad r:=\sqrt{x_r^2+y_r^2},\quad
\psi(r):=\begin{cases}
5r,&r<0.2,\\
2-5r,&0.2\le r<0.4,\\
0,&r\ge0.4.
\end{cases}
\end{equation*}

We take the CFL number $K_{\rm CFL}=0.475$ and compute the numerical solution until the final time $t=1$ on a uniform $128\times128$ mesh
for $\ve=10^{-\alpha}$ with $\alpha=1,\dots,6$, and report the obtained local Mach number, defined as $\|\bm u\|_2/\sqrt{\gamma}$, in Figure
\ref{fig43} along with its initial distribution. According to what is expected due to the AP feature of the scheme, the shape of the vortex
is preserved and no evident dependency on $\ve$ can be observed.
\begin{figure}[ht!]
\centerline{\includegraphics[width=0.33\textwidth]{gresho_initial}}
\vskip8pt
\centerline{\includegraphics[width=0.33\textwidth]{gresho_T1_eps0.1_CFL0.475}\hspace*{0.1cm}
\includegraphics[width=0.33\textwidth]{gresho_T1_eps0.01_CFL0.475}\hspace*{0.1cm}
\includegraphics[width=0.33\textwidth]{gresho_T1_eps0.001_CFL0.475}}
\vskip8pt
\centerline{\includegraphics[width=0.33\textwidth]{gresho_T1_eps0.0001_CFL0.475}\hspace*{0.1cm}
            \includegraphics[width=0.33\textwidth]{gresho_T1_eps1e-05_CFL0.475}\hspace*{0.1cm}
	    \includegraphics[width=0.33\textwidth]{gresho_T1_eps1e-06_CFL0.475}}
\caption{\sf Example 2: Initial local Mach number independently of $\ve$ (top) and local Mach number at $t=1$ for different values of $\ve$.
\label{fig43}}
\end{figure}

\subsubsection*{Example 3---Baroclinic Vorticity Generation}
In this example taken from \cite{noelle2014weakly}, we consider a low-Mach-number flow with $\ve=0.05$ involving an acoustic wave, which
moves within two density layers in the computational domain $[-\frac{1}{\ve},\frac{1}{\ve}]\times[0,\frac{2}{5\ve}]$ subject to the periodic
boundary conditions. The acoustic wave induces different accelerations in the two density layers, which results in rotational excitation and
in the formation of a long-wavelength sinusoidal shear layer. Due to the interaction with the acoustic wave, such a shear layer becomes
unstable, and several Kelvin-Helmholtz-type unstable structures originate from it.

The initial conditions are
\begin{equation*}
\begin{aligned}
&\rho(x,y,0)=1+\frac{\ve}{2000}[1+\cos(\ve\pi x)]+4.5\ve y-\begin{cases}0,&0\le y\le\frac{1}{5\ve},\\1.8,&\mbox{otherwise,}\end{cases}\\
&u(x,y,0)=\frac{\sqrt{\gamma}}{2}[1+\cos(\ve\pi x)],\quad v(x,y,0)\equiv0,\quad p(x,y,0)=1+\frac{\ve\gamma}{2}[1+\cos(\ve\pi x)].
\end{aligned}
\end{equation*}

The numerical solution is computed with the CFL number $K_{\rm CFL}=0.475$ until the final time $t=20$ on a $800\times160$ uniform mesh. The
density at times $t=0$, $10$, and $20$ is plotted in Figure \ref{fig44}. Since the solution develops instabilities, no convergence is
expected in this example; see \cite{zeifang2020novel}. One can, however, observe that the underlying physics is correctly captured.
\begin{figure}[ht!]
\centerline{\includegraphics[width=0.9\textwidth]{baroclinic_WC_SOLUTION_OCT_0000000.0000000.dat_eps0.05_CFL0.475}}
\centerline{\includegraphics[width=0.9\textwidth]{baroclinic_WC_SOLUTION_OCT_0000010.0000000.dat_eps0.05_CFL0.475}}
\centerline{\includegraphics[width=0.9\textwidth]{baroclinic_WC_SOLUTION_OCT_0000020.0000000.dat_eps0.05_CFL0.475}}
\caption{\sf Example 3: Density at different times.\label{fig44}}
\end{figure}

\subsubsection*{Example 4---Double Shear Layer Problem}
In the following test case, originally introduced in \cite{bell1989second} for the incompressible Navier–Stokes equations and subsequently 
adopted in, e.g., \cite{zeifang2020novel,boscarino2018all,weinan1994numerical} in the context of compressible Euler equations in the
low-Mach-number regime, a shear layer develops and the AP property of the proposed scheme can be assessed. In particular, we would like to
check whether the scheme maintains its consistency for small values of $\ve$, that is, in the almost incompressible regime.

The initial conditions,
\begin{align*}
&\rho(x,y,0)\equiv\frac{\pi}{15},\quad
u(x,y,0)=\left\{\begin{aligned}
&\tanh\!\left[15\Big(\frac{y}{\pi}-\frac{1}{2}\Big)\right],&&y\le\pi,\\[0.5ex]
&\tanh\!\left[15\Big(\frac{3}{2}-\frac{y}{\pi}\Big)\right],&&\mbox{otherwise},
\end{aligned}\right.\quad 
v(x,y,0)=0.05\sin x,\\
&p(x,y,0)\equiv\frac{1}{\gamma}.
\end{align*}
are prescribed in the computational domain $[0,2\pi]\times[0,2\pi]$ subject to the periodic boundary conditions. The initial vorticity
$\omega:=v_x-u_y$, where the derivatives are approximated using the second-order central differences, is plotted in Figure \ref{fig45}.
\begin{figure}[ht!]
\centerline{\includegraphics[width=0.33\textwidth]{double_shear_layer_initial_vorticity}}
\caption{\sf Example 4: Initial vorticity.\label{fig45}}
\end{figure}

We compute the numerical solutions for $\ve=10^{-\alpha}$ with $\alpha=1,\dots,6$ until the final time $t=10$ on a $256\times256$ uniform
mesh using $K_{\rm CFL}=0.1$. Figures \ref{fig46} and \ref{fig47} display the vorticity at times $t=6$ and $t=10$, respectively, for
different $\ve$. The obtained results are consistent with those reported in \cite{zeifang2020novel,boscarino2018all}. Moreover, no
macroscopic dependence on $\ve$ is observed, providing further evidence of the AP property of the proposed DF-FV scheme.
\begin{figure}[ht!]
\centerline{\includegraphics[width=0.33\textwidth]{double_shear_layer_WC_SOLUTION_OCT_0000006.0000000.dat_eps0.1_CFL0.1}\hspace*{0.1cm}
            \includegraphics[width=0.33\textwidth]{double_shear_layer_WC_SOLUTION_OCT_0000006.0000000.dat_eps0.01_CFL0.1}\hspace*{0.1cm}
            \includegraphics[width=0.33\textwidth]{double_shear_layer_WC_SOLUTION_OCT_0000006.0000000.dat_eps0.001_CFL0.1}}
\vskip8pt
\centerline{\includegraphics[width=0.33\textwidth]{double_shear_layer_WC_SOLUTION_OCT_0000006.0000000.dat_eps0.0001_CFL0.1}\hspace*{0.1cm}
            \includegraphics[width=0.33\textwidth]{double_shear_layer_WC_SOLUTION_OCT_0000006.0000000.dat_eps1e-05_CFL0.1}\hspace*{0.1cm}
            \includegraphics[width=0.33\textwidth]{double_shear_layer_WC_SOLUTION_OCT_0000006.0000000.dat_eps1e-06_CFL0.1}}
\caption{\sf Example 4: Vorticity at $t=6$ for different values of $\ve$. $K_{\rm CFL}=0.1$.\label{fig46}}
\end{figure}
\begin{figure}[ht!]
\centerline{\includegraphics[width=0.33\textwidth]{double_shear_layer_WC_SOLUTION_OCT_0000010.0000000.dat_eps0.1_CFL0.1}\hspace*{0.1cm}
            \includegraphics[width=0.33\textwidth]{double_shear_layer_WC_SOLUTION_OCT_0000010.0000000.dat_eps0.01_CFL0.1}\hspace*{0.1cm}
            \includegraphics[width=0.33\textwidth]{double_shear_layer_WC_SOLUTION_OCT_0000010.0000000.dat_eps0.001_CFL0.1}}
\vskip8pt
\centerline{\includegraphics[width=0.33\textwidth]{double_shear_layer_WC_SOLUTION_OCT_0000010.0000000.dat_eps0.0001_CFL0.1}\hspace*{0.1cm}
            \includegraphics[width=0.33\textwidth]{double_shear_layer_WC_SOLUTION_OCT_0000010.0000000.dat_eps1e-05_CFL0.1}\hspace*{0.1cm}
            \includegraphics[width=0.33\textwidth]{double_shear_layer_WC_SOLUTION_OCT_0000010.0000000.dat_eps1e-06_CFL0.1}}
\caption{\sf Example 4: The same as in Figure \ref{fig46}, but at $t=10$.\label{fig47}}
\end{figure}

We remark that the simulations remain stable for larger values of the CFL number. However, the use of larger $K_{\rm CFL}$ may lead to a
noticeable increase in the amount of the numerical diffusion for very small values of $\ve\lesssim10^{-3}$. To illustrate this, we
re-compute the solution with $K_{\rm CFL}=0.475$ for $\ve=10^{-4}$, $10^{-5}$, and $10^{-6}$ and plot the obtained results (for $t=6$) in
Figure \ref{fig48}. As one can clearly see, the numerical solution is now substantially more diffusive compared with those reported in the
bottom row of Figure \ref{fig46}.
\begin{figure}[ht!]
\centerline{\includegraphics[width=0.33\textwidth]{double_shear_layer_WC_SOLUTION_OCT_0000006.0000000.dat_eps0.0001_CFL0.475}\hspace*{0.1cm}
            \includegraphics[width=0.33\textwidth]{double_shear_layer_WC_SOLUTION_OCT_0000006.0000000.dat_eps1e-05_CFL0.475}\hspace*{0.1cm}
            \includegraphics[width=0.33\textwidth]{double_shear_layer_WC_SOLUTION_OCT_0000006.0000000.dat_eps1e-06_CFL0.475}}
\caption{\sf Example 4: The same as in Figure \ref{fig46}, but for $K_{\rm CFL}=0.475$ (left).\label{fig48}}
\end{figure}

\subsubsection*{Example 5---Explosion Problem}
In the last numerical example, we consider an explosion problem taken from \cite{ToroBook}. The initial data,
$$
(\rho,u,v,p)(x,y,0)=\begin{cases}(1,0,0,1),&\sqrt{x^2+y^2}<0.4,\\(0.125,0,0,0.1),&\mbox{otherwise},\end{cases}
$$
are prescribed in the computational domain $[-1,1]\times[-1,1]$ subject to the free boundary conditions.

The main objective of this test is to verify that the proposed AP DF-FV scheme remains accurate and stable also in the high-Mach-number
regime, in which strong shocks and contact discontinuities may be present. To this end, we perform simulations for several values of $\ve$.
For $\ve=1$, $0.9$, $0.6$, and $0.3$ the final times are $t=0.25$, $0.2$, $0.15$, and $0.08$, respectively. The surface plots of the density
$\rho$ computed on a uniform mesh with $400\times400$ cells using $K_{\rm CFL}=0.475$ are reported in Figure \ref{fig49}, where one can see
that the obtained solutions oscillation-free and their nonsmooth features are accurately resolved for all values of $\ve$.
\newcommand{\trimfig}[1]{\includegraphics[width=0.35\textwidth, trim=140 260 190 250, clip]{#1}}
\begin{figure}[ht!]
\centerline{\trimfig{new_sod2d1.0_rho_3d_time_scheme_IMEX_DeC2_Prim_Explicit_Cons_CFL0.475}\hspace*{0.5cm}
\trimfig{new_sod2d0.9_rho_3d_time_scheme_IMEX_DeC2_Prim_Explicit_Cons_CFL0.475}}
\vskip8pt
\centerline{\trimfig{new_sod2d0.6_rho_3d_time_scheme_IMEX_DeC2_Prim_Explicit_Cons_CFL0.475}\hspace*{0.5cm}
\trimfig{new_sod2d0.3_with_timestep_relaxation_rho_3d_time_scheme_IMEX_DeC2_Prim_Explicit_Cons_CFL0.475}}
\caption{\sf Example 5: Surface plot of density for different values of $\ve$ and corresponding times.\label{fig49}}
\end{figure}

To further assess the correctness of the computed solutions, we plot their one-dimensional (1-D) slices along the diagonal $y=x$ in Figure
\ref{fig410} together with the corresponding slices of the reference solution, which was obtained using the second-order semi-discrete CU
scheme from \cite{KLin} on a much finer mesh with $2000\times2000$ cells using the CFL number $0.2$ and the three-stage third-order strong
stability preserving (SSP) Runge-Kutta method \cite{GKS,GST}. As one can clearly see, the computed solutions shows a perfect agreement with
the reference ones, and the discontinuities locations are correctly captured.
\begin{figure}[ht!]
\centerline{\includegraphics[width=0.8\textwidth]{new_sod2d1.0_timescheme_IMEX_DeC2_Prim_Explicit_Cons_CFL0.475_V_compare_diag}}
\vskip3pt
\centerline{\includegraphics[width=0.8\textwidth]{new_sod2d0.9_timescheme_IMEX_DeC2_Prim_Explicit_Cons_CFL0.475_V_compare_diag}}
\vskip3pt
\centerline{\includegraphics[width=0.8\textwidth]{new_sod2d0.6_timescheme_IMEX_DeC2_Prim_Explicit_Cons_CFL0.475_V_compare_diag}}
\vskip3pt
\centerline{\includegraphics[width=0.8\textwidth]{new_sod2d0.3_with_timestep_relaxation_timescheme_IMEX_DeC2_Prim_Explicit_Cons_CFL0.475_V_compare_diag}}
\caption{\sf Example 5: 1-D slices of the computed solutions along $y=x$ for different values of $\ve$ and at different times: $\ve=1$,
$t=0.25$ (top row), $\ve=0.9$, $t=0.2$ (second row), $\ve=0.6$, $t=0.15$ (third row), and $\ve=0.3$, $t=0.08$ (bottom row).\label{fig410}}
\end{figure}

\begin{remark}
We stress that in this example, both the denominators in \eref{4.48f} will vanish if $\rho_{\max}$ and $p_{\max}$ are computed using
\eref{4.3}. The modification \eref{4.9} will prevent this from happening, but the resulting time steps might still be too big to guarantee
the stability of the AP DF-FV method. Indeed, an adjustment to the initial time steps is needed for the small values of $\ve$=$0.6$ or
$0.3$. We therefore set $\dt=10^{-4}$ for the first 10 time steps.

In fact, a similar stability problem occurs even when explicit methods are used as a direct evaluation of the wave speeds based on the
available data leads to an underestimate of the real wave speeds; see \cite{toro2020bounds} and \cite[Remark 6]{micalizzitoro2024}. This is
particularly true for Riemann problems with zero initial velocities. Instead of artificially selecting very small few first time steps, one
may compute the time steps based on the rigorous wave speed bounds. Another less efficient alternative would be to reduce $K_{\rm CFL}$ to
$0.2$ also $\ve=0.6$ and $0.05$ also $\ve=0.3$. This will lead to the computed solutions, which are practically indistinguishable from those
reported in Figures \ref{fig49} and \ref{fig410}.
\end{remark}
\begin{remark}
Let us remark that discontinuities are unlikely to occur in low-Mach-number flows. Consequently, the above tests with $\ve=0.6$ and $0.3$
should be regarded as rather ``academic'' and are primiry intended to demonstrate that the proposed AP DF-FV scheme is capable of handling
discontinuities even in the low-Mach-number regime.
\end{remark}

\section{Conclusion}\label{sec5}
We have presented a novel asymptotic-preserving (AP) numerical method for the compressible Euler equations that is effective across all
Mach-number regimes, including the low-Mach-number one where standard explicit schemes become inefficient. The key idea is a new hyperbolic
splitting, inspired by the flux-splitting approach introduced in \cite{haack2012all}. The new splitting is applied to a primitive
(nonconservative) formulation of the Euler equations, which enables one to design an efficient semi-implicit (SI) time discretizations. Our
splitting isolates stiff linear terms, which are discretized semi-implicitly: this leads to a well-posed linear elliptic problem, which
helps ensure the AP property of the resulting scheme.
	
To overcome the well-known difficulties associated with the use of nonconservative formulations in the presence of discontinuities, we
implement the proposed AP scheme within the recently introduced dual formulation framework. In this approach, the conservative and primitive
systems are solved simultaneously, and their resulting solutions are post-processed to ensure the correct capturing of discontinuities whileretaining the AP property of the primitive-based SI approach.
	
The proposed AP dual formulation finite-volume (DF-FV) method has been thoroughly validated on several benchmarks ranging from the fully
compressible to the nearly incompressible regime, demonstrating both high accuracy and robustness of the method. Future work will focus on
extending the AP DF-FV framework to more complex systems and on developing higher-order spatial and temporary discretizations.
	
\begin{acknowledgment}
The work of L. Micalizzi was supported in part by the LeRoy B. Martin, Jr. Distinguished Professorship Foundation. The work of A. Kurganov
was supported in part by NSFC grant W2431004. The work of A. Chertock was supported in part by NSF grant DMS-2208438. 
\end{acknowledgment}

\bibliography{ref}
\bibliographystyle{siam}
\end{document}
